% !TeX encoding = UTF-8
% !TeX program = pdflatex
% IEEEtran Journal Template -- Migrated from Elsevier cja-template.tex
\documentclass[lettersize,journal]{IEEEtran}

%% ---- Packages (only those compatible with IEEEtran) ----
\usepackage{amsmath,amsfonts,amssymb}
\usepackage{algorithm}
\usepackage{algorithmic}    % IEEEtran recommends algorithmic (not algorithmicx)
\usepackage{array}
\usepackage{bm}
\usepackage{booktabs}
\usepackage{multirow}
\usepackage{graphicx}
\usepackage{url}
\usepackage{cite}           % IEEE-style numeric citations
\usepackage{textcomp}
\usepackage{stfloats}       % allows figure* at bottom
% \usepackage[caption=false,font=normalsize,labelfont=sf,textfont=sf]{subfig}
\usepackage[caption=false,font=footnotesize,labelfont=sf,textfont=sf]{subfig}
\usepackage{siunitx}

% ---- Highlight/revision markup (see new_highlight/RULES.md) ----
\usepackage{xcolor}
\usepackage[normalem]{ulem}   % normalem: keep \emph as italics, not underline
% Added text: plain blue, NO wavy underline / underline
\newcommand{\hladd}[1]{{\color{blue}#1}}
% Deleted inline text: red with a strike through the middle of the text
\newcommand{\hldel}[1]{{\color{red}\sout{#1}}}
% Deleted display equation: red, with a strike through the centre of the formula
\newcommand{\hldeleq}[1]{%
  {\color{red}%
   \sbox0{$\displaystyle #1$}%
   \rlap{\usebox0}%
   \rule[0.5\dimexpr\ht0-\dp0\relax]{\wd0}{0.6pt}}}
% Graphics live in the original manuscript folder
\graphicspath{{../orign_paper/}{../HIL_fig/}{../Re2_6/}{../Re3_8/}{../Re3_9/}}
% 3D (V10) result figures for the revised simulation section
\newcommand{\vfig}[1]{../../new_sim_fig/fig_and_data/figures_v10/#1}

% theorem/lemma environments
\newtheorem{theorem}{Theorem}
\newtheorem{lemma}{Lemma}

\hyphenation{op-tical net-works semi-conduc-tor IEEE-Xplore}

% ----------------------------------------------------------------
\begin{document}

% ---- Title ----
% \title{Multi-Agent Reinforcement Learning-Based Cooperative Interception Strategy
%   for Maneuvering Swarm Attack Targets with Considering Target Assignment}

% \title{Trust-Aware Multi-Agent Reinforcement Learning-Based Cooperative Interception Strategy for Maneuvering Swarm Attack Targets Considering Rapid Consensus Target Assignment}
\title{Trust-Aware Multi-Agent Reinforcement Learning-Based Cooperative Interception Strategy for Maneuvering Swarms with Target Assignment}

% ---- Authors ----
% \author{Shuyi~Gao,~Defu~Lin,~Duo~Zheng,~and~Zhichen~Chu%
% \thanks{S. Gao, D. Lin, D. Zheng, and Z. Chu are with the School of Aerospace
% Engineering, Beijing Institute of Technology, Beijing 100081, China
% (e-mail: zhengduohello@126.com).}%
% \thanks{Manuscript received 2025.}}

%----双盲投稿--
\author{}
\markboth{Under review}%
{Anonymous Submission}

% ---- Running heads ----
\markboth{IEEE Transactions on Automation Science and Engineering}%
{Gao \MakeLowercase{\textit{et al.}}: MARL-Based Cooperative Interception for Maneuvering Swarm Targets}

\maketitle

% ---- Abstract ----
\begin{abstract}
The swarm attack of unmanned aerial vehicles (UAVs) poses a significant threat to future
battlefield operations.
Achieving efficient and precise cooperative interception against attacking UAV swarms with
highly dynamic maneuvering capabilities and adaptive penetration strategies poses a critical
challenge for UAV interceptor swarms.
This study addresses the cooperative interception of maneuvering targets in swarm attacks by
decomposing the problem into two key tasks: target assignment and cooperative interception
guidance.
For target assignment, an Improved Dung Beetle Optimization (IDBO) algorithm is proposed to
dynamically balance exploration and exploitation, achieving highly efficient and rapid
assignments.
% For cooperative interception guidance, a multi-agent deep reinforcement learning-based method
% is developed to achieve precise and time-coordinated interception of adversarial UAV swarms
% employing advanced penetration strategies.
% For cooperative interception guidance, a multi-agent deep reinforcement learning framework—augmented by kinematics-aware attention mechanisms and trust-modulated exploration—is developed to achieve precise and time-coordinated interception of adversarial UAV swarms employing advanced penetration strategies.
For cooperative interception guidance, an Attention-Residual and Trust-aware Multi-Agent Proximal Policy Optimization (ART-MAPPO) framework is developed to achieve precise and time-coordinated interception of adversarial UAV swarms employing advanced penetration strategies.
The simulation results validate the effectiveness of the proposed intelligent cooperative
interception method, demonstrating superior convergence speed, robustness, and significantly
improved interception success rates in swarm UAV cooperative interception scenarios.
\end{abstract}

{\linespread{0.9}\selectfont
\hspace*{0.0em}\textbf{\textit{Note to Practitioners—}}\textbf{This paper is motivated by practical UAV swarm defense scenarios, where multiple interceptors must cooperatively engage maneuvering aerial targets under limited communication and onboard maneuver constraints. In such applications, target assignment and guidance are closely coupled, and conventional centralized approaches may be difficult to implement efficiently as the engagement scale increases. To address this issue, this paper develops a two-stage cooperative interception framework that combines distributed target assignment with a trust-aware multi-agent reinforcement learning guidance method. The proposed approach is intended to improve coordination consistency among interceptors assigned to the same target and to enhance robustness against target maneuvers while considering maneuver saturation and command smoothness. This work may be useful for practical large-scale multi-UAV interception and other cooperative autonomous systems requiring scalable distributed decision-making and coordinated guidance.}\par\vspace{0.5\baselineskip}
}

% ---- Keywords ----
\begin{IEEEkeywords}
Cooperative interception, deep reinforcement learning, centralized training with decentralized
execution, target assignment, UAV swarm.
\end{IEEEkeywords}


%% ===========================================================
\section{Introduction}
\label{sec:introduction}
\IEEEPARstart{A}{s} warfare evolves toward intelligent swarm operations, the likelihood of
encountering attacks from adversarial UAV swarms significantly increases during combat.
Consequently, the efficient interception of incoming swarm targets becomes a critical
imperative for future battlefield operations.
To address the challenge of cooperative interception against maneuvering attacking swarms with
adaptive penetration strategies, there is an urgent need to advance academic research in swarm
target interception technologies.
This study aims to enhance interceptors with intelligent cooperative interception capabilities,
enabling effective navigation of complex challenges in modern battlefield environments.

Multi-target cooperative interception encompasses two key tasks: target assignment and
cooperative interception.
The effectiveness of multi-target interception hinges on the optimality of target assignment,
establishing it as a primary focus in swarm target interception strategies.
As the number of UAVs increases, the complexity of target assignment escalates significantly,
resulting in a multi-constrained, multi-parameter optimization problem.
The objective is to develop a rapid and efficient target assignment method for allocating
specific attacking targets to interceptors in operational scenarios~\cite{KLINE2019226}.
The target assignment problem involves complex constraints, and the non-dominated sorting
algorithm can dynamically generate a set of reference points with good convergence and
distribution performance~\cite{WANG2025952,bai2024hybrid,zou2024moea}.
Moreover, employing an online operator selection mechanism can enhance the efficiency of the
algorithm~\cite{ahmad2024fuzzy,ming2024constrained}.
The target assignment problem can be solved using multi-information particle swarm optimization
to find the solution that yields the highest gain~\cite{mishra2015weapon,wei2020particle}.
Introducing information from surrounding particles can enhance the stability of the traditional
particle swarm algorithm while also addressing the issue of local convergence in the target
assignment problem~\cite{long2024improved}.
The genetic algorithm is a common approach for solving the target assignment
problem~\cite{he4914655modelling,li2020multi,sun2022multi,chang2018solving}, and a GA with a
novel crossover operator can further enhance algorithmic efficiency~\cite{weng2024diversity}.
Ant colony optimization is also an effective strategy for solving the target assignment
problem~\cite{li2024comprehensive}, and its combination with a greedy algorithm can improve
the accuracy of parallel computation~\cite{yi2024solving}.
Under multiple constraints, researchers have proposed various heuristic algorithms to solve the
target assignment problem, such as neural networks~\cite{li2024comprehensive}, auction
algorithm~\cite{wang2024collaborative}, tabu search~\cite{yao2024tabu}, simulated annealing
algorithm~\cite{khurshid2024hybrid}, VLSN~\cite{yue2024graph}.
However, in cooperative interception scenarios, the aforementioned algorithms struggle to
handle the highly dynamic nature of the battlefield due to their high time complexity, and they
fail to meet the dual requirements of both speed and optimality.

The effectiveness of the terminal guidance phase in multi-target interception is significantly
influenced by the cooperative guidance law, a critical factor requiring thorough consideration.
The target assignment strategy transforms swarm operations into a coordinated multi-to-one
interception problem, with current research focusing on three key areas: angular coordination,
temporal coordination, and intelligent cooperation.
Significant progress has been made in temporal coordination.
Classical proportional navigation-based guidance laws enable synchronized multi-missile attacks
on a single target~\cite{tao2024analytical}.
Adaptive control methods adjust attack timing based on real-time operational conditions, making
them suitable for intercepting maneuvering targets~\cite{wang2025generalized,zhang2022multiple}.
Sliding Mode Control (SMC) is widely adopted for cooperative
interception~\cite{sinha2020supertwisting}, with SMC incorporating miss distance and attack
time constraints enhancing interception
accuracy~\cite{shen2025three,jin2025cooperative,hu2019sliding}.
Considerable research has also been conducted on angular coordination.
From a traditional control perspective, backstepping control~\cite{li2024adaptive} and
leader-follower cooperative methods~\cite{he2017impact} significantly improve interception
efficiency.
SMC with angular constraints is prevalent in cooperative interception, encompassing approaches
based on impact angle constraints~\cite{wang2025hybrid} and line-of-sight angle
constraints~\cite{li2024near}.
Optimal control methods applied to impact angle-constrained cooperative guidance enhance
precision by integrating impact angle and line-of-sight angle
constraints~\cite{wang2025adaptive,tao2023optimal}.

Recent advancements in intelligent technologies have driven extensive research into cooperative
interception strategies for UAV swarms, with intelligent algorithms demonstrating substantial
potential in addressing complex multi-target scenarios~\cite{chen2024enhanced,wu2020collaborative}.
Optimization-based approaches, such as Particle Swarm Optimization (PSO), excel in
collaborative attack problems by autonomously optimizing resource allocation and maximizing
operational efficiency~\cite{zhao2024hierarchical,wu2020collaborative}.
Search-based algorithms, like A*~\cite{liao2024research}, enable efficient path planning for
coordinated attacks.
Control-based methods, including Model Predictive Control
(MPC)~\cite{kim2024adaptive}, provide robust guidance under dynamic constraints.
Moreover, the integration of neural networks with reinforcement learning, particularly deep
reinforcement learning (DRL), enhances the adaptability and intelligence of interception
strategies~\cite{luo2022uav,zhan2022multiple,wang2021uav}.
DRL enables UAV swarms to learn complex policies from high-dimensional data, optimizing target
assignment and maneuvering in real-time to counter maneuvering targets~\cite{sun2024memory}.
In contrast, conventional cooperative interception guidance methods rely heavily on
linearization and simplifying assumptions to design feasible guidance laws~\cite{zhao2025global}.
These methods often struggle to address the complexity of dynamic, non-linear environments
involving maneuvering swarm targets.
Data-driven intelligent cooperative interception methods, by contrast, leverage extensive
datasets to learn adaptive strategies, significantly reducing dependence on restrictive
assumptions~\cite{fainkich2025cooperative}.
By modeling complex interactions through neural networks and DRL, these methods achieve greater
robustness and flexibility.
Despite advancements, intelligent cooperative interception methods struggle with maneuvering
swarm targets in complex dynamic environments due to limited model generalization, low training
efficiency, and low interception success rates, particularly in large-scale swarm combat
requiring precise temporal coordination~\cite{horyna2023decentralized}.
Addressing these challenges requires further research into scalable, generalizable, and
efficient algorithms to fully realize the potential of intelligent cooperative interception.

In summary, existing cooperative interception methods for countering attacking UAV swarms with
adaptive penetration capabilities suffer from low interception success rates in complex dynamic
environments.
Specifically, target assignment methods exhibit low allocation efficiency, while cooperative
interception guidance strategies struggle with terminal overload saturation and insufficient
adaptability to dynamic target maneuvers in practical scenarios, hindering effective swarm
coordination.
Consequently, more research is necessary to address these challenges.

This study proposes a swarm cooperative interception method based on multi-agent reinforcement
learning and IDBO algorithm, enhancing target assignment efficiency and guidance strategy
performance to achieve precise and efficient interception of maneuvering targets.
The proposed intelligent cooperative interception algorithm offers the following advantages.

\begin{enumerate}
\item A distributed consensus-based target assignment scheme is proposed, which integrates an
improved Dung Beetle Optimization (IDBO) algorithm with a combat-dominance assessment model.
This design supports rapid many-to-one allocation under limited communication, leading to more
efficient and reliable target assignment for cooperative interception.

\item ART-MAPPO is proposed as a CTDE-based cooperative interception guidance learning
framework, which employs an attention--residual feature extractor and a GRU temporal encoder,
and integrates trust-modulated exploration with a dual-clip PPO update and adaptive KL
regularization. This design yields faster convergence and improved training efficiency in
highly maneuverable swarm-adversarial environments.

\item \hldel{A cooperative interception method augmented with a reinforcement learning-based reward
function is introduced, enabling precise temporal coordination.} \hladd{While reward shaping and
clipped policy updates are standard reinforcement-learning mechanisms, the contribution here lies
in instantiating them specifically for the maneuvering-swarm engagement: the time-to-go consensus
reward is formulated as a learned impact-time (salvo) constraint that defeats sequential, piecemeal
interception, and the dual-clip update is cast as a load-factor saturation guard bounding the
commanded terminal-overload increment under aggressive exploration.} This strategy enhances
interception success rates against large-scale maneuvering swarm targets\hldel{ by addressing dynamic
coordination challenges}\hladd{ while preserving terminal-load safety}.
\end{enumerate}

The paper is structured as follows: Section~\ref{sec:method} presents relevant problem modeling and foundational theory. Section~\ref{sec:target_assignment} and Section~\ref{sec:cooperative_algorithm} provide a detailed description of our approach, including its framework, modules, algorithms, and rationale. Section~\ref{sec:simulation} covers experiments, results, and discussions. Lastly, Section~\ref{sec:conclusion} summarizes the paper and outlines future directions.

%% ===========================================================
\section{Backgrounds and Preliminaries}
\label{sec:method}

In this section, we first introduce the modeling of multi-target interception of UAVs, and
then review reinforcement learning methods.

\subsection{Problem Modeling}

The multi-target intelligent interception problem can be decomposed into two subproblems:
target assignment and intelligent cooperative interception, with the swarm combat scenario
illustrated in Fig.~\ref{fenpei}.
In this scenario, the attacking swarm aims to penetrate a high-value target area, while the
defending swarm seeks to fully intercept the incoming UAVs.
During the target assignment phase, the defending swarm assesses its own flight status and
that of the attacking swarm using onboard sensors, subsequently allocating each defender to an
appropriate interception target via the target assignment strategy.
In the intelligent cooperative interception phase, defenders assigned to the same target employ
coordinated maneuvering strategies to intercept and neutralize the target.

\begin{figure}[!t]
\centering
\includegraphics[width=2.8in]{fenpei.pdf}
\caption{Cluster target interception problem.}
\label{fenpei}
\end{figure}

The object of study in this paper is the fixed-wing UAV\hldel{, and the interception problem is
reduced to a two-dimensional horizontal plane as shown in}\hladd{.}~Fig.~\ref{jianmo} \hldel{without loss of
generality}\hladd{shows the horizontal projection used to define the heading and line-of-sight
(LOS) azimuth, whereas the guidance model and simulations propagate the complete
three-dimensional position and velocity}.

\begin{figure}[!t]
\centering
\includegraphics[width=2.6in]{pic02.pdf}
\caption{Schematic diagram of confrontation and interception.}
\label{jianmo}
\end{figure}

The \hladd{three-dimensional point-mass} kinematic model of a single UAV \hldel{can be simplified as} \hladd{is}
\begin{equation}
\label{yundongxue}
\hldeleq{\left\{\begin{matrix}\dot{x} =V\cos\gamma \\
\dot{y} =V\sin\gamma \end{matrix}\right.}
\end{equation}
\begin{equation}
\label{yundongxue3d}
{\color{blue}\left\{\begin{aligned}
\dot{x}&=V\cos\theta\cos\gamma,\\
\dot{y}&=V\cos\theta\sin\gamma,\\
\dot{z}&=V\sin\theta,
\end{aligned}\right.}
\end{equation}
where \hldel{\emph{x}, \emph{y} is the position coordinates of the UAV in two-dimensional space,
\emph{V} is the velocity magnitude of the UAV, and $\gamma$ is the navigation angle.}%
\hladd{$(x,y,z)$ are the three-dimensional position coordinates, $V$ is the velocity magnitude,
$\gamma$ is the heading angle, and $\theta$ is the flight-path (pitch) angle.}

In an interception mission, \hldel{the relative motion relationship between the attacking and
defending UAVs can be described as:}\hladd{let $\mathbf{p}_A,\mathbf{p}_D\in\mathbb{R}^3$ and
$\mathbf{v}_A,\mathbf{v}_D\in\mathbb{R}^3$ denote the attacker and defender states. Their
three-dimensional relative geometry is}
\begin{equation}
\label{gongfang}
\hldeleq{\left\{\begin{matrix}\dot{r} =V_{A}\cos(\gamma_{A}-q)-V_{D}\cos(\gamma_{D}-q) \\
r\dot{q} =V_{A}\sin(\gamma_{A}-q)-V_{D}\sin(\gamma_{D}-q)\end{matrix}\right.}
\end{equation}
\begin{equation}
\label{gongfang3d}
{\color{blue}\begin{aligned}
\mathbf{r}_{AD}&=\mathbf{p}_A-\mathbf{p}_D,\qquad
r=\|\mathbf{r}_{AD}\|_2,\\
\boldsymbol{\ell}_{AD}&=\mathbf{r}_{AD}/r,\qquad
\dot{r}=\boldsymbol{\ell}_{AD}^{\top}(\mathbf{v}_A-\mathbf{v}_D).
\end{aligned}}
\end{equation}
\hldel{where $r$ is the inter-UAV distance, $q$ is the line-of-sight angle, $V_{A}$ ($V_{D}$) is the
velocity of the attacking (defending) UAV, and $\gamma_{A}$ ($\gamma_{D}$) is the
corresponding heading angle.}%
\hladd{Here, $r$ is the three-dimensional inter-UAV distance and $\boldsymbol{\ell}_{AD}$ is the
LOS unit vector; its horizontal projection provides the plan-view geometry shown in
Fig.~\ref{jianmo}, in which $q$ denotes the LOS azimuth and $V_{A}$ ($V_{D}$),
$\gamma_{A}$ ($\gamma_{D}$) are the attacker (defender) speed and heading.}

\hldel{In the interception phase, the relationship between UAV flight state and acceleration is:}%
\hladd{The relationship between the flight state and the three commanded load channels is}
\begin{equation}
\label{guozai}
\hldeleq{\left\{\begin{matrix}\dot{V} =n_{x} g \\
\dot{\gamma} =n_{y}g/V \end{matrix}\right.}
\end{equation}
\begin{equation}
\label{guozai3d}
{\color{blue}\left\{\begin{aligned}
\dot{V}&=n_x g,\\
\dot{\gamma}&=\frac{n_y g}{V\cos\theta},\\
\dot{\theta}&=\frac{n_z g}{V},
\end{aligned}\right.}
\end{equation}
where $n_{x}$ is the axial \hldel{acceleration command, $n_{y}$ is the normal acceleration command,}%
\hladd{load-factor command, $n_y$ and $n_z$ are the horizontal and vertical normal load-factor
commands, respectively,} and $g$ is the gravitational acceleration.
\hladd{Thus, $n_y$ primarily changes heading, while $n_z$ changes the flight-path angle and
altitude.}
In the multi-agent setting of Section~\ref{sec:cooperative_algorithm}, these variables carry
an agent subscript $i$.

During cooperative interception, the remaining impact time is defined as:
\begin{equation}
\label{shijian}
t_{go}=\frac{r}{\left|\dot{r}\right|}
\end{equation}

The \hladd{three-dimensional} Zero-Effort Miss (ZEM) is defined as\hldel{:}
\begin{equation}
\label{ZEM}
\hldeleq{ZEM = \frac{r^{2}\cdot\dot{q}}{\left|\dot{r}\right|}}
\end{equation}
\begin{equation}
\label{ZEM3d}
{\color{blue}\operatorname{ZEM}=\left\|\mathbf{r}_{AD}+(\mathbf{v}_A-\mathbf{v}_D)\,t_{go}\right\|_2.}
\end{equation}

\subsection{Reinforcement Learning}

Reinforcement learning (RL) is built on the Markov Decision Process (MDP), comprising state
space $\mathcal{S}$, action space $\mathcal{A}$, reward function $R$, policy $\pi$, and state
transition function $P$.
At each step, the agent selects action $a$ from state $s$, receives reward $r$, and
transitions to $s'$.
The key value functions are:
\begin{equation}
\label{RLV}
V^{\pi}(s)=E_{\pi}\left[\sum_{t=0}^{\infty}\gamma_{\mathrm{RL}}^{t}r_{t}\mid s_{0}=s\right]
\end{equation}
\begin{equation}
\label{RLQ}
Q^{\pi}(s,a)=E_{\pi}\left[\sum_{t=0}^{\infty}\gamma_{\mathrm{RL}}^{t}r_{t}\mid s_{0}=s,a_{0}=a\right]
\end{equation}
where $\gamma_{\mathrm{RL}}\in[0,1]$ is the discount factor.






\section{Target Assignment via Distributed Consensus-Based IDBO Algorithm}
\label{sec:target_assignment}

This section formulates the cooperative target assignment as a weighted bipartite matching problem and solves it via an Improved Dung Beetle Optimization (IDBO) algorithm embedded in a distributed consensus framework.

\subsection{Problem Formulation}
\label{subsec:problem_formulation}

Let $\mathcal{U} = \{u_1, u_2, \dots, u_M\}$ and $\mathcal{T} = \{t_1, t_2, \dots, t_N\}$ denote the interceptor and target sets, respectively. Figure \ref{km} illustrates the bipartite assignment structure.

\begin{figure}[!t]
\centering
\includegraphics[width=2.5in]{KM1.pdf}
\caption{Schematic diagram of weighted dichotomy.}
\label{km}
\end{figure}

The binary decision variable $X_{ij}=1$ if UAV $u_i$ is assigned to target $t_j$ and 0 otherwise. The edge weight $P_{ij}$ quantifies the interception probability:

\begin{equation}
\label{eq:interception_probability}
P_{ij} = w_1 \cdot P_{\text{ZEM}}(i,j) + w_2 \cdot P_{\sigma}(i,j)
\end{equation}
where $w_1 + w_2 = 1$, $w_1, w_2 \geq 0$. The Zero-Effort Miss (ZEM) component is:

\begin{equation}
\label{eq:zem_prob}
P_{\text{ZEM}}(i,j) = \exp\left[-\frac{1}{2}\left(\frac{\text{ZEM}(i,j)}{\overline{\text{ZEM}}}\right)^2\right]
\end{equation}
where $\overline{\text{ZEM}} = \frac{1}{MN}\sum_{i=1}^{M}\sum_{j=1}^{N} \text{ZEM}(i,j)$.

The angular advantage component is:

\begin{equation}
\label{eq:angular_prob}
P_{\sigma}(i,j) = \exp\left[-\frac{1}{2}\left(\frac{S_{\sigma}(i,j)}{\overline{S_{\sigma}}}\right)^2\right]
\end{equation}
where $\overline{S_{\sigma}} = \frac{1}{MN}\sum_{i=1}^{M}\sum_{j=1}^{N} S_{\sigma}(i,j)$, with $S_{\sigma}(i,j) = |\gamma_{ij} - q_{ij}|$ denoting the heading error. The objective minimizes the expected number of surviving targets. Under many-to-one assignment, the joint survival probability of target $t_j$ is the product of individual failure probabilities:

\begin{equation}
\label{eq:optimization_formulation}
\begin{aligned}
\min_{\mathbf{X}} \quad & J = \sum_{j=1}^{N} \left[ \prod_{i=1}^{M} \left(1 - P_{ij}\right)^{X_{ij}} \right] \\
\text{s.t.} \quad & \sum_{j=1}^{N} X_{ij} = 1, \quad \forall i \in \{1,\dots,M\} \\
& \sum_{i=1}^{M} X_{ij} \geq 1, \quad \forall j \in \{1,\dots,N\} \\
& \sum_{i=1}^{M} X_{ij} \leq L_{\max}, \quad \forall j \in \{1,\dots,N\} \\
& X_{ij} \in \{0,1\}, \quad \forall i,j
\end{aligned}
\end{equation}
where $\prod_{i=1}^{M} (1 - P_{ij})^{X_{ij}}$ yields the survival probability of $t_j$, which decreases geometrically with the number of assigned interceptors.

\hladd{A sensitivity study on the weighting pair $(w_1,w_2)$ clarifies how tactical priority reshapes the assignment. Sweeping $w_1$ from $0$ to $1$ over $20$ defenders versus $8$ maneuvering attackers (12 randomized geometries) shows a clear trade-off: as $w_1$ increases toward the ZEM term, the mean selected miss drops by about $56.9\%$ (from $1152$\,m to $676$\,m), whereas the mean terminal aspect error grows from roughly $61.5^{\circ}$ to $93.5^{\circ}$; giving more weight to $P_{\sigma}$ reverses this, minimizing aspect error at the cost of reachability. Crucially, the aggregate objective is robust to this choice: the expected number of surviving targets $J$ stays below $0.18$ across the entire sweep (i.e., above $97.7\%$ expected neutralization), and the weakest-target coverage probability varies by under $3\%$. Consequently, $(w_1,w_2)$ acts as a tactical dial---weighting reachability versus firing-geometry quality---rather than an effectiveness switch, and the near-balanced setting $w_1 \approx 0.55$ used here yields a favorable compromise between the two objectives.}

\subsection{Distributed Consensus-Based IDBO Framework}
\label{subsec:idbo_algorithm}

The IDBO algorithm solves \eqref{eq:optimization_formulation} by alternating between local metaheuristic optimization and distributed consensus-based conflict resolution. Let $\mathbf{X}^{(k)} = [x_{ij}^{(k)}] \in \{0,1\}^{M \times N}$ denote the assignment matrix at iteration $k$, where $x_{ij}^{(k)}=1$ indicates UAV $u_i$ is assigned to target $t_j$. The integrated optimization is:

\begin{equation}
\label{eq:hierarchical}
\mathcal{A}_{\text{IDBO}}^{(k+1)} = \mathcal{G}_{\text{consensus}} \circ \mathcal{F}_{\text{optimizer}} \left( \mathbf{X}^{(k)}, \boldsymbol{\Theta}^{(k)} \right)
\end{equation}
where $\mathcal{F}_{\text{optimizer}}$ denotes the improved DBO operator and $\mathcal{G}_{\text{consensus}}$ the consensus-based conflict resolver. The parameter set $\boldsymbol{\Theta}^{(k)} = \{\mathbf{A}^{(k)}, \mathbf{P}^{(k)}, \mathbf{V}^{(k)}\}$ aggregates the adversarial advantage, interception probability, and velocity matrices at iteration $k$.

Each UAV $u_i$ maintains a preference vector $\mathbf{x}_i \in [0,1]^{N}$, where $x_{ij}$ encodes the normalized affinity for target $t_j$. The fitness function balances interception efficiency against resource saturation:

\begin{equation}
\label{eq:idbo_fitness}
\begin{split}
\mathcal{F}_i(\mathbf{x}_i) =& \sum_{j=1}^{N} \sigma(x_{ij}) \cdot P_{ij} + \lambda_1 \sum_{j=1}^{N} \sigma(x_{ij}) \cdot A_{ij} \\
&- \lambda_2 \sum_{j=1}^{N} \Psi_{ij}(\mathbf{x}_i, \hat{\mathbf{X}}_{\mathcal{N}_i})
\end{split}
\end{equation}
where the three terms correspond to interception probability, adversarial advantage ($\lambda_1$-weighted), and an over-saturation penalty ($\lambda_2$-weighted). The penalty $\Psi_{ij}$ is:

\begin{equation}
\label{eq:saturation_penalty}
\begin{split}
\Psi_{ij}(\mathbf{x}_i, \hat{\mathbf{X}}_{\mathcal{N}_i}) =& \sigma(x_{ij}) \cdot \max\Biggl(0, \Biggl(\sum_{k \in \mathcal{N}_i, k \neq i} \sigma(\hat{x}_{kj}) \\
&+ \sigma(x_{ij})\Biggr) - L_{\max}\Biggr)
\end{split}
\end{equation}
where $\mathcal{N}_i$ is the neighbor set of $u_i$, $\hat{x}_{kj}$ the estimated preference of neighbor $u_k$, and $L_{\max}$ the per-target interceptor capacity. The penalty activates only when the estimated assignment count exceeds $L_{\max}$, steering the swarm toward capacity-feasible solutions.

The adversarial advantage $A_{ij}$ fuses temporal, energetic, and geometric factors:

\begin{equation}
\label{eq:adversarial_advantage}
\begin{split}
A_{ij} =& \exp\left(-\frac{\Delta t_{ij}}{T_{\max}}\right) \cdot \left(1 - \frac{\|\mathbf{v}_i - \mathbf{v}_j\|_2}{v_{\max}}\right) \\
&\cdot \max\left(0, \cos(\theta_{ij})\right) - \sum_{k\neq i} \frac{P_{kj}}{P_{ij} + P_{kj}} \cdot \sigma(\hat{x}_{kj})
\end{split}
\end{equation}
where $\Delta t_{ij}$ is the estimated time-to-intercept; $\mathbf{p}_i=[x_i, y_i]^\top$ and $\mathbf{v}_i=[\dot{x}_i, \dot{y}_i]^\top$ denote the position and velocity vectors of interceptor $u_i$; $\mathbf{p}_j=[x_j, y_j]^\top$ and $\mathbf{v}_j=[\dot{x}_j, \dot{y}_j]^\top$ denote the position and velocity vectors of target $t_j$; $\theta_{ij}$ is the aspect angle associated with the engagement geometry between $u_i$ and $t_j$; and $v_{\max}$ is the maximum reference speed used for normalization. The three multiplicative factors encode temporal, energetic, and geometric advantage, respectively, while the subtracted term penalizes competition from other interceptors with high preferences for the same target.

The IDBO update mechanism enhances the standard DBO through four adversarial-information-driven operations:

\begin{align}
\text{Rolling:} \quad & \mathbf{x}_i^{(t+1)} = \mathbf{x}_i^{(t)} + \alpha \cdot \nabla_{\mathbf{x}_i} \mathcal{F}_i(\mathbf{x}_i^{(t)}) \circ \mathbf{A}_i \nonumber \\
& \quad + \beta \cdot \mathbf{r}_i \cdot \exp\left(-\frac{t}{T}\right) \label{eq:rolling_update} \\
\text{Dancing:} \quad & \mathbf{x}_i^{(t+1)} = \mathbf{x}_i^{(t)} \odot \exp\left(\gamma \cdot \tanh(\mathbf{A}_i - \bar{A}_i)\right) \nonumber \\
& \quad + \mathbf{n}_i, \quad \mathbf{n}_i \sim \mathcal{N}(0, \sigma_d^2 \mathbf{I}) \label{eq:dancing_update} \\
\text{Breeding:} \quad & \mathbf{x}_i^{(t+1)} = \mathbf{x}_i^{(t)} + \delta \cdot \left(\mathbf{x}_{\text{best}}^{(t)} - \mathbf{x}_i^{(t)}\right) \odot \mathbf{A}_i \nonumber \\
& \quad + \delta' \cdot \left(\mathbf{x}_{\text{second}}^{(t)} - \mathbf{x}_i^{(t)}\right) \label{eq:breeding_update} \\
\text{Stealing:} \quad & \mathbf{x}_i^{(t+1)} = \mathbf{x}_i^{(t)} \nonumber \\
& \quad + \eta \cdot \sum_{k \in \mathcal{K}_i} \frac{A_{ij}}{\sum_{l \in \mathcal{K}_i} A_{lj}} \cdot \left(\mathbf{x}_k^{(t)} - \mathbf{x}_i^{(t)}\right) \label{eq:stealing_update}
\end{align}
where $\mathbf{A}_i = [A_{i1}, A_{i2}, \dots, A_{iN}]^\top$, $\bar{A}_i$ is the mean of $\mathbf{A}_i$, $\mathbf{r}_i$ and $\mathbf{n}_i$ are random vectors, $\mathcal{K}_i$ represents UAVs competing for the same targets as $u_i$, and $\alpha, \beta, \gamma, \delta, \delta', \eta$ are adaptive coefficients that decay linearly with iteration $t$. \hladd{The coefficients $\alpha,\beta,\gamma,\delta,\delta',\eta$ together with the dancing noise scale $\sigma_d$ follow a common linear schedule $c_t=c_0\,(1-t/T)$, where $T$ is the IDBO iteration count, while the elite-recombination weights remain bounded. This schedule is what governs convergence stability. The stochastic parts of the operators---the rolling probe $\beta\,\mathbf{r}_i\exp(-t/T)$ and the dancing perturbation $\mathbf{n}_i\sim\mathcal{N}(0,\sigma_d^2\mathbf{I})$---scale with $c_t$ and vanish as $c_t\!\to\!0$, so the update degenerates into a deterministic, advantage-weighted local refinement followed by the advantage-adaptive binarization step. A step sequence whose stochastic amplitude decays to zero while retaining an ascent direction satisfies a Robbins--Monro-type sufficient condition~\cite{robbins1951stochastic} for settling at a stable fixed point rather than orbiting it, which is precisely the behavior measured by the swarm-disagreement metric $\Gamma^{(k)}$ in \eqref{eq:convergence}. A controlled ablation that freezes the coefficients confirms this mechanism: non-decaying coefficients keep injecting perturbations and leave a persistent late-stage oscillation in the working assignment, whereas the linear schedule contracts the swarm onto a stable, low-cost equilibrium.} Continuous preferences are binarized via advantage-adaptive thresholding:

\begin{equation}
\label{eq:thresholding}
X_{ij} = 
\begin{cases}
1, & \text{if } \sigma(x_{ij}) > \tau \cdot \left(1 + \frac{A_{ij}}{\max_k A_{ik}}\right) \\
0, & \text{otherwise}
\end{cases}
\end{equation}
where $\tau \in (0,1)$ is a base threshold. Each UAV maintains a bundle $\mathcal{B}_i$, bid vector $\mathbf{c}_i = [c_{i1}, \dots, c_{iN}]$, winner list $\mathbf{z}_i = [z_{i1}, \dots, z_{iN}]$, and timestamp vector $\mathbf{y}_i = [y_{i1}, \dots, y_{iN}]$. The bid for target $t_j$ governs the distributed auction:

\begin{equation}
\label{eq:bid_formulation}
c_{ij} = \mathcal{F}_i(\mathbf{x}_i) \cdot \sigma(x_{ij}) \cdot \exp\left(\frac{A_{ij} - \bar{A}_i}{\sigma_A}\right)
\end{equation}
where $\bar{A}_i$ and $\sigma_A$ are the mean and standard deviation of $u_i$'s adversarial advantages. The fitness--preference product captures local assignment quality, while the exponential term amplifies bids for targets where $u_i$ holds above-average advantage. To support cooperative saturation attacks, the consensus mechanism employs top-$K$ bidding. Each UAV maintains a winner list $\mathcal{Z}_{j} = \{(u_k, c_{kj})\}$ for $t_j$, containing up to $L_{\max}$ highest bids:

\begin{equation}
\label{eq:cbba_update}
\mathcal{Z}_{j}^{(k+1)} = \text{Top}_{L_{\max}} \left( \mathcal{Z}_{j}^{(k)} \cup \bigcup_{l \in \mathcal{N}_i} \mathcal{Z}_{j,l}^{(k)} \cup \{(u_i, c_{ij}^{(k+1)})\} \right)
\end{equation}
where $\text{Top}_{L_{\max}}(\cdot)$ retains the $L_{\max}$ highest-bid entries. UAV $u_i$ secures a slot for $t_j$ if $|\mathcal{Z}_{j}| < L_{\max}$ or its bid exceeds $\min_{(u_k, c_{kj}) \in \mathcal{Z}_{j}} c_{kj}$, thereby displacing the weakest contributor and prioritizing the most capable interceptors.

As depicted in the procedural flowchart (Fig.~\ref{fig:idbo_flowchart}), the algorithm iterates over three interleaved phases—local optimization (Eqs.~\ref{eq:rolling_update}--\ref{eq:stealing_update}), consensus resolution (Eqs.~\ref{eq:thresholding}--\ref{eq:cbba_update}), and advantage update—until convergence:

\begin{align}
\text{Phase 1: } & \mathbf{x}_i^{(k+1)} = \mathcal{F}_{\text{optimizer}}\left(\mathbf{x}_i^{(k)}, \mathcal{F}_i, \mathbf{A}_i^{(k)}, \mathbf{P}^{(k)}\right) \label{eq:phase1} \\
\text{Phase 2: } & \{\mathbf{z}_i^{(k+1)}, \mathcal{B}_i^{(k+1)}\} = \nonumber \\ 
& \mathcal{G}_{\text{consensus}}\left(\{\sigma(\mathbf{x}_i^{(k+1)})\}, \{\mathbf{c}_i^{(k+1)}\}, \{\mathbf{y}_i^{(k)}\}\right) \label{eq:phase2} \\
\text{Phase 3: } & \mathbf{A}_i^{(k+1)} = \mathcal{H}\left(\mathbf{z}_i^{(k+1)}, \mathcal{B}_i^{(k+1)}, \mathbf{V}_i^{(k+1)}, \mathbf{V}_T^{(k+1)}\right) \label{eq:phase3}
\end{align}

\begin{figure}[!t]
\centering
\includegraphics[width=3.2in]{assignment.pdf}
\caption{Flowchart of the distributed consensus-based IDBO algorithm, illustrating the algorithmic iteration flow across local metaheuristic optimization, auction-based consensus, and adversarial information updates.}
\label{fig:idbo_flowchart}
\end{figure}

Here, $\mathcal{H}$ updates adversarial advantages given the consensus assignment and velocity vectors $\mathbf{V}_i$, $\mathbf{V}_T$. Following the iteration flow in Fig.~\ref{fig:idbo_flowchart}, convergence is declared when the swarm-wide assignment disagreement falls below $\epsilon$:

\begin{equation}
\label{eq:convergence}
\Gamma^{(k)} = \frac{1}{M} \sum_{i=1}^{M} \left\| \mathbf{z}_i^{(k)} - \frac{1}{|\mathcal{N}_i|} \sum_{j \in \mathcal{N}_i} \mathbf{z}_j^{(k)} \right\|_0 < \epsilon
\end{equation}
where $\|\cdot\|_0$ counts disagreeing entries and $\epsilon$ is a convergence threshold. \hldel{The algorithm guarantees monotonic fitness improvement: $\Phi^{(k+1)} \geq \Phi^{(k)}$ a.s., where $\Phi^{(k)} = \sum_{i=1}^{M} \mathcal{F}_i(\mathbf{x}_i^{(k)})$. Under a connected communication graph with finite message delays, consensus is reached in at most $N \cdot D$ iterations ($D$: graph diameter). The converged assignment $\mathbf{X}^*$ constitutes an $\epsilon$-approximate Nash equilibrium: $\mathcal{F}_i(\mathbf{x}_i^*) \geq \mathcal{F}_i(\mathbf{x}_i') - \epsilon$, $\forall \mathbf{x}_i' \in \mathcal{X}_i$, $\forall i$.} \hladd{Three properties characterize this iteration under a connected communication graph with bounded message delay and elitist retention of best-so-far solutions. First, the \emph{incumbent} (best-so-far) swarm fitness $\Phi_{\mathrm{best}}^{(k)}=\max_{k'\le k}\sum_{i=1}^{M}\mathcal{F}_i(\mathbf{x}_i^{(k')})$ is non-decreasing by construction, and because the decaying schedule $c_t=c_0(1-t/T)$ drives the operators' stochastic amplitude to zero while retaining an ascent direction, the working assignment settles at a fixed point rather than orbiting it---the standard Robbins--Monro condition for stochastic-approximation convergence~\cite{robbins1951stochastic}. Second, the auction-based consensus phase~\eqref{eq:bid_formulation}--\eqref{eq:cbba_update} inherits the guarantees of the consensus-based bundle algorithm~\cite{choi2009consensus}: with a diminishing-marginal-gain (submodular) score it reaches a conflict-free assignment in at most $N\!\cdot\!D$ communication rounds ($D$: graph diameter). Third, the resulting fixed point $\mathbf{X}^\ast$ is $\epsilon$-\emph{stable}---no interceptor can raise its own fitness by more than $\epsilon$ through a unilateral reassignment, $\mathcal{F}_i(\mathbf{x}_i^\ast)\ge\mathcal{F}_i(\mathbf{x}_i')-\epsilon$ for all $\mathbf{x}_i'\in\mathcal{X}_i$ and all $i$, which follows directly from the auction best-response at the converged bid prices~\cite{choi2009consensus}. The empirical convergence of the swarm-disagreement metric $\Gamma^{(k)}$ and the near-linear growth of consensus rounds with $D$ are reported in the supplementary response material. The complete per-interceptor procedure is summarized in Algorithm~\ref{alg:idbo}, where each step is cross-referenced to the corresponding equation for reproducibility.}

\begin{algorithm}[!t]
\color{blue}
\caption{Distributed Consensus-Based IDBO for Cooperative Target Assignment (per interceptor $u_i$)}
\label{alg:idbo}
\begin{algorithmic}[1]
\REQUIRE Interceptor set $\mathcal{U}$, target set $\mathcal{T}$, probability matrix $\mathbf{P}$~\eqref{eq:interception_probability}, neighbors $\mathcal{N}_i$, capacity $L_{\max}$, IDBO iterations $T$, population size $N_{\mathrm{pop}}$, threshold $\epsilon$
\ENSURE Converged assignment $\mathbf{X}^\ast$
\STATE Initialize preference population $\{\mathbf{x}_i\}$, winner lists $\mathcal{Z}_j\!\leftarrow\!\emptyset$, timestamps $\mathbf{y}_i$, iteration $k\!\leftarrow\!0$
\REPEAT
  \STATE {\itshape // Phase 1: local metaheuristic optimization}
  \FOR{$t=1$ to $T$}
    \STATE Set decaying coefficients $c_t=c_0(1-t/T)$ for $\{\alpha,\beta,\gamma,\delta,\delta',\eta,\sigma_d\}$
    \STATE Evaluate fitness $\mathcal{F}_i$~\eqref{eq:idbo_fitness} and advantage $\mathbf{A}_i$~\eqref{eq:adversarial_advantage}
    \STATE Update $\mathbf{x}_i$ by Rolling/Dancing/Breeding/Stealing~\eqref{eq:rolling_update}--\eqref{eq:stealing_update}
  \ENDFOR
  \STATE Binarize preferences by advantage-adaptive thresholding~\eqref{eq:thresholding}
  \STATE {\itshape // Phase 2: auction-based distributed consensus}
  \STATE Compute bids $c_{ij}$~\eqref{eq:bid_formulation}
  \STATE Exchange winner lists with $\mathcal{N}_i$; update top-$L_{\max}$ set $\mathcal{Z}_j$~\eqref{eq:cbba_update}
  \STATE {\itshape // Phase 3: adversarial-information update}
  \STATE Refresh advantages $\mathbf{A}_i^{(k+1)}$ from the consensus assignment
  \STATE Compute swarm disagreement $\Gamma^{(k)}$~\eqref{eq:convergence};\; $k\!\leftarrow\!k+1$
\UNTIL{$\Gamma^{(k)}<\epsilon$}
\STATE \textbf{return} consensus assignment $\mathbf{X}^\ast$
\end{algorithmic}
\end{algorithm}

The per-UAV per-iteration complexity decomposes as:

\begin{equation}
\label{eq:complexity}
C_{\text{total}} = O(T \cdot P \cdot N^2) + O(N \cdot |\mathcal{N}_i|) + O(N)
\end{equation}
where $T$ is the IDBO iteration count, $P = O(N)$ the population size, and $|\mathcal{N}_i|$ the neighborhood size. The three terms account for optimization, consensus communication, and advantage update, respectively. Communication cost is $O(N \cdot |\mathcal{N}_i|)$ and memory is $O(P \cdot N + N)$ per UAV. \hladd{Because the target set faced by each interceptor and its neighborhood size $|\mathcal{N}_i|$ are fixed by the engagement and the communication topology rather than by the swarm size $M$, the per-UAV per-iteration cost above is independent of $M$; the total per-round cost of the swarm therefore grows only linearly with $M$. Scalability is governed not by this arithmetic cost but by how fast the distributed consensus in \eqref{eq:convergence} settles. Under a connected graph with a per-hop communication latency $T_{\mathrm{comm}}$, each target's winning bids propagate across the network in at most $D$ hops ($D$: graph diameter), so the wall-clock time to reach consensus scales as $O(N\!\cdot\!D\!\cdot\!T_{\mathrm{comm}})$. Empirically this time stays nearly constant as $M$ grows at fixed diameter, and increases only linearly with the communication delay and with $D$; a bounded delay thus slows, but does not prevent, convergence of the distributed consensus. Detailed complexity, scalability, and communication-delay measurements are reported in the supplementary response material.}













%% ===========================================================
\section{ART-MAPPO: Attention-Residual and Trust-Aware MAPPO for Cooperative Swarm Interception}
\label{sec:cooperative_algorithm}

Building on the distributed consensus-based IDBO target assignment, we propose the Attention-Residual and Trust-aware MAPPO (ART-MAPPO) algorithm for the cooperative guidance of defending UAV swarms under the Centralized Training with Decentralized Execution (CTDE) paradigm. ART-MAPPO augments the standard MAPPO framework with three synergistic modules: a kinematics-aware attention-residual network for extracting cross-coupled engagement features, a trust-modulated exploration mechanism for discovering cooperative tactics within the safe flight envelope, and a dual-clip PPO with adaptive KL penalty for constraining guidance command variations against actuator saturation.

\subsection{Enhanced Situational Awareness via Attention-Residual Network}
\label{subsec:attention_residual}

The highly dynamic nature of swarm interception requires interceptors to process complex, coupled kinematic states. Let $\mathbf{o}_i^{(t)} \in \mathbb{R}^{n_o}$ denote the observation vector for interceptor $i$ at time step $t$, as defined in Section~\ref{subsec:state_reward}. A base feature encoder $f_{\mathrm{enc}}$, comprising a two-layer MLP with hidden dimension $d_{\mathrm{enc}}$ and Layer Normalization, projects the raw observation into a high-dimensional latent space:
\begin{equation}
\label{eq:encoder}
\begin{split}
\mathbf{h}^{(0)} =& \operatorname{LN}\Bigl(\operatorname{ReLU}\bigl(\mathbf{o}_i^{(t)}\mathbf{W}_{\mathrm{enc}}^{(1)} + \mathbf{b}_{\mathrm{enc}}^{(1)}\bigr) \\
&\cdot \mathbf{W}_{\mathrm{enc}}^{(2)} + \mathbf{b}_{\mathrm{enc}}^{(2)}\Bigr) \in \mathbb{R}^{d}
\end{split}
\end{equation}
where $\mathbf{W}_{\mathrm{enc}}^{(1)} \in \mathbb{R}^{n_o \times d_{\mathrm{enc}}}$ and $\mathbf{W}_{\mathrm{enc}}^{(2)} \in \mathbb{R}^{d_{\mathrm{enc}} \times d}$ are learnable projection matrices.

\emph{1) Kinematics-Aware Multi-Head Attention:} To capture the nonlinear spatio-temporal coupling among engagement variables (e.g., the interdependence between LOS rate $\dot{q}$ and time-to-go synchronization), we reshape $\mathbf{h}^{(0)}$ into a sequence of $n_o$ feature tokens $\mathbf{H}^{(0)} \in \mathbb{R}^{n_o \times (d/n_o)}$ and apply multi-head self-attention. For each head $j \in \{1, \dots, H\}$, the query, key, and value projections are:
\begin{equation}
\label{eq:qkv}
\mathbf{Q}_j = \mathbf{H}^{(0)}\mathbf{W}_j^{Q}, \quad \mathbf{K}_j = \mathbf{H}^{(0)}\mathbf{W}_j^{K}, \quad \mathbf{V}_j = \mathbf{H}^{(0)}\mathbf{W}_j^{V}
\end{equation}
where $\mathbf{W}_j^{Q}, \mathbf{W}_j^{K} \in \mathbb{R}^{(d/n_o) \times d_k}$ and $\mathbf{W}_j^{V} \in \mathbb{R}^{(d/n_o) \times d_v}$, with $d_k = d_v = d/H$. The scaled dot-product attention for head $j$, incorporating an observability mask $\mathbf{M} \in \mathbb{R}^{n_o \times n_o}$, is:
\begin{align}
\mathbf{A}_j &= \mathrm{softmax}\!\left(\frac{\mathbf{Q}_j\mathbf{K}_j^\top}{\sqrt{d_k}} + \mathbf{M}\right) \in \mathbb{R}^{n_o \times n_o} \label{eq:attention} \\
\mathrm{head}_j &= \mathbf{A}_j\mathbf{V}_j \in \mathbb{R}^{n_o \times d_v} \nonumber
\end{align}
where $M_{u,v} = -\infty$ isolates unobservable kinematic channels (e.g., LOS rate during sensor blanking) and $M_{u,v} = 0$ permits information flow. The multi-head outputs are concatenated, projected, and added back via a residual shortcut:
\begin{equation}
\label{eq:mha_out}
\begin{split}
\mathbf{h}_{\mathrm{attn}} =& \mathbf{h}^{(0)} + \mathrm{Dropout}_{p}\Bigl(\operatorname{Flatten}\bigl( \\
&\left[\mathrm{head}_1 \| \cdots \| \mathrm{head}_H\right]\bigr) \mathbf{W}^{O}\Bigr)
\end{split}
\end{equation}
where $\mathbf{W}^{O} \in \mathbb{R}^{d \times d}$ is the output projection and $\operatorname{Flatten}(\cdot)$ reshapes the token sequence back to $\mathbb{R}^d$. The residual connection ensures that primitive flight-state components propagate losslessly to downstream layers.

\emph{2) Residual Block with Layer Normalization:} To deepen the nonlinear representational capacity without gradient vanishing---a critical issue when extreme overload commands produce large-magnitude gradients---a stack of $L$ residual blocks follows the attention module. For the $\ell$-th block ($\ell \in \{1, \dots, L\}$) with input $\mathbf{h}^{(\ell)}$ (where $\mathbf{h}^{(1)} = \mathbf{h}_{\mathrm{attn}}$), the forward pass is:
\begin{align}
\mathbf{z}_1^{(\ell)} &= \operatorname{ReLU} \left( \operatorname{LN}_1^{(\ell)} \left( \mathbf{h}^{(\ell)}\mathbf{W}_1^{(\ell)} + \mathbf{b}_1^{(\ell)} \right) \right) \label{eq:res_block_inner1} \\
\mathbf{z}_2^{(\ell)} &= \operatorname{LN}_2^{(\ell)} \left( \mathbf{z}_1^{(\ell)}\mathbf{W}_2^{(\ell)} + \mathbf{b}_2^{(\ell)} \right) \label{eq:res_block_inner2} \\
\mathbf{h}^{(\ell+1)} &= \operatorname{ReLU} \left( \mathbf{z}_2^{(\ell)} + \mathbf{h}^{(\ell)} \right) \label{eq:res_block_out}
\end{align}
where $\mathbf{W}_1^{(\ell)}, \mathbf{W}_2^{(\ell)} \in \mathbb{R}^{d \times d}$ are orthogonally initialized. The identity shortcut yields the Jacobian:
\begin{equation}
\label{eq:jacobian}
\begin{split}
\frac{\partial \mathbf{h}^{(\ell+1)}}{\partial \mathbf{h}^{(\ell)}} =& \operatorname{diag}\!\left(\mathbb{I}\!\left[\mathbf{z}_2^{(\ell)} + \mathbf{h}^{(\ell)} > \mathbf{0}\right]\right) \\
&\cdot \left(\mathbf{I} + \frac{\partial \mathbf{z}_2^{(\ell)}}{\partial \mathbf{h}^{(\ell)}}\right)
\end{split}
\end{equation}
Since the ReLU gate preserves the identity component on active units, gradient norms satisfy $\left\|\frac{\partial \mathbf{h}^{(\ell+1)}}{\partial \mathbf{h}^{(\ell)}}\right\| \geq \left\|\operatorname{diag}(\mathbb{I}[\cdot])\right\| > 0$, preventing vanishing gradients across $L$ blocks.

\emph{3) End-to-End Guidance Command Generation:} The refined feature $\mathbf{h}^{(L+1)}$ is fed into a GRU cell to integrate temporal engagement history:
\begin{equation}
\label{eq:gru}
\mathbf{c}_t = \operatorname{GRU}\!\left(\mathbf{h}^{(L+1)},\, \mathbf{c}_{t-1};\, \mathbf{W}_{\mathrm{GRU}}\right), \quad \mathbf{c}_t \in \mathbb{R}^{d_{\mathrm{rnn}}}
\end{equation}
where $\mathbf{W}_{\mathrm{GRU}}$ denotes the GRU parameters and $\mathbf{c}_{t-1}$ is the hidden state from the previous guidance step. The actor network then parameterizes a diagonal Gaussian over the continuous overload action space:
\begin{equation}
\label{eq:actor_dist}
\begin{split}
\pi_{\theta_i}\!\left(a_i \mid \mathbf{o}_i^{(t)}\right) &= \mathcal{N}\!\left(\boldsymbol{\mu}_{\theta_i}(\mathbf{c}_t),\; \operatorname{diag}\!\left(\boldsymbol{\sigma}_{\theta_i}^2(\mathbf{c}_t)\right)\right) \\
a_i &= [n_{x_i},\, n_{y_i}]^\top \in \mathcal{A}
\end{split}
\end{equation}
where $\boldsymbol{\mu}_{\theta_i}, \boldsymbol{\sigma}_{\theta_i}^2: \mathbb{R}^{d_{\mathrm{rnn}}} \to \mathbb{R}^2$ are linear output heads. The action is clipped to the feasible envelope $\mathcal{A}$ after sampling. Concurrently, the centralized critic shares the same attention-residual backbone but receives the concatenated global state:
\begin{equation}
\label{eq:critic_input}
s_t = \left[\mathbf{o}_1^{(t)} \| \mathbf{o}_2^{(t)} \| \cdots \| \mathbf{o}_N^{(t)}\right] \in \mathbb{R}^{N \cdot n_o}
\end{equation}
to produce the state-value estimate $V_\phi(s_t) \in \mathbb{R}$, enabling centralized advantage computation during training while each actor executes independently from local observations $\mathbf{o}_i^{(t)}$ at deployment.

\subsection{Trust-Modulated Tactical Exploration}
\label{subsec:trust_exploration}

Based on the high-fidelity features extracted by the attention-residual backbone, we introduce a trust-based mechanism to dynamically modulate the exploration policy. This mechanism mathematically balances the exploitation of known guidance laws against the discovery of complex cooperative tactics (e.g., flanking or simultaneous impact).

\emph{1) Trust Evolution:} Let $R_i^{(k)} = \sum_{t=0}^{T} \gamma_{\mathrm{RL}}^t r_i^{(t)}$ be the discounted cumulative return of interceptor $i$ in episode $k$, where $\gamma_{\mathrm{RL}} \in (0,1)$ is the RL discount factor (distinct from the heading angle $\gamma_{D_i}$) and $T$ is the episode horizon. The swarm-level running statistics are maintained as exponential moving averages:
\begin{align}
\mu_R &\leftarrow (1-\rho)\,\mu_R + \rho\,\frac{1}{N}\sum_{i=1}^{N} R_i^{(k)} \label{eq:running_stats_mu} \\
\sigma_R^2 &\leftarrow (1-\rho)\,\sigma_R^2 + \rho\,\frac{1}{N}\sum_{i=1}^{N}\!\left(R_i^{(k)} - \mu_R\right)^2 \label{eq:running_stats_sigma}
\end{align}
where $\rho \in (0,1)$ is the statistics update rate. The normalized relative performance of interceptor $i$ is then:
\begin{equation}
\label{eq:trust_normalized}
\tilde{R}_i^{(k)} = \frac{R_i^{(k)} - \mu_R}{\sqrt{\sigma_R^2 + \epsilon_0}}
\end{equation}
where $\epsilon_0 > 0$ is a numerical stabilizing constant. The trust level $\mathcal{T}_i^{(k)} \in (0,1)$ evolves via first-order exponential smoothing with a sigmoid activation:
\begin{equation}
\label{eq:trust_evolution}
\mathcal{T}_i^{(k+1)} = (1-\alpha_T)\,\mathcal{T}_i^{(k)} + \alpha_T\,\sigma\!\left(\tau_T \tilde{R}_i^{(k)}\right)
\end{equation}
where $\sigma(x) = 1/(1+e^{-x})$, $\alpha_T \in (0,1)$ is the trust update rate and $\tau_T > 0$ is the temperature parameter controlling the sensitivity of trust to performance deviations. \hldel{When $R_i^{(k)} > \mu_R$, the sigmoid output exceeds 0.5, driving $\mathcal{T}_i$ upward and reducing exploration; conversely, below-average performance lowers $\mathcal{T}_i$ and amplifies exploration.} \hladd{Equation~\eqref{eq:trust_evolution} is a first-order exponential-smoothing (contraction) update, not a monotonic accumulator: at each episode it relaxes $\mathcal{T}_i^{(k)}$ toward the instantaneous target $\sigma(\tau_T\tilde{R}_i^{(k)})\in(0,1)$ at rate $\alpha_T$. Because $\alpha_T\in(0,1)$ and this target is bounded in $(0,1)$, the recursion is a contraction of modulus $(1-\alpha_T)$; hence $\mathcal{T}_i^{(k)}$ never leaves $(0,1)$ and converges toward the moving equilibrium $\mathcal{T}_i^{\star}=\sigma(\tau_T\tilde{R}_i^{(k)})$ set by the current relative performance. Consistently above-average performance ($R_i^{(k)}>\mu_R$) holds this target above $0.5$, so trust settles at and \emph{maintains} a high equilibrium---suppressing exploration---rather than growing without bound, while persistently below-average performance pulls the equilibrium below $0.5$ and amplifies exploration. This convergent-smoothing behavior, not a monotonic increment, is precisely what sustains the stable high-trust regime.}

\emph{2) Blended Exploration Policy:} The final action for interceptor $i$ is drawn from a two-component mixture: with probability $1-\beta(\mathcal{T}_i)$ from the learned policy, and with probability $\beta(\mathcal{T}_i)$ from a structured heuristic distribution, i.e.,
\begin{equation}
\label{eq:exploration_policy}
a_i \sim \begin{cases} \pi_{\theta_i}(\cdot \mid \mathbf{o}_i^{(t)}) & \text{with prob. } 1 - \beta(\mathcal{T}_i) \\ \pi_{\mathrm{guided}}^{i}(\cdot \mid \mathbf{o}_i^{(t)}) & \text{with prob. } \beta(\mathcal{T}_i) \end{cases}
\end{equation}
where the modulation factor $\beta(\mathcal{T}_i) = 1 - \mathcal{T}_i \in (0,1)$ is monotonically decreasing in trust. The heuristic distribution $\pi_{\mathrm{guided}}^{i}$ is a discrete mixture over three tactically informative actions within the feasible command space $\mathcal{A} = [n_{x}^{\min}, n_{x}^{\max}] \times [n_{y}^{\min}, n_{y}^{\max}]$:
\begin{equation}
\label{eq:guided}
\begin{split}
\pi_{\mathrm{guided}}^{i}(a_i \mid \mathbf{o}_i^{(t)}) =& \omega_1\,\delta(a_i - a_{\mathrm{opt}}) + \omega_2\,\delta(a_i - a_{\mathrm{probe}}) \\
&+ \omega_3\,\mathcal{U}(\mathcal{A}), \quad \sum_{k=1}^3\omega_k=1
\end{split}
\end{equation}
Here, $a_{\mathrm{opt}} = \arg\max_{a \in \mathcal{A}} Q_\phi(\mathbf{o}_i^{(t)}, a)$ mimics the classical proportional navigation command for exploitation, $a_{\mathrm{probe}} = \arg\min_{a \in \mathcal{A}} Q_\phi(\mathbf{o}_i^{(t)}, a)$ deliberately executes extreme boundary maneuvers to probe the aerodynamic envelope, and $\mathcal{U}(\mathcal{A})$ maintains stochastic coverage of the remaining action space. The expected action under the full mixture policy is:
\begin{equation}
\label{eq:expected_action}
\begin{split}
\mathbb{E}[a_i] =& \bigl(1-\beta(\mathcal{T}_i)\bigr)\,\boldsymbol{\mu}_{\theta_i}(\mathbf{c}_t) \\
&+ \beta(\mathcal{T}_i)\,\bigl(\omega_1\,a_{\mathrm{opt}} + \omega_2\,a_{\mathrm{probe}} + \omega_3\,\bar{a}_{\mathcal{A}}\bigr)
\end{split}
\end{equation}
where $\bar{a}_{\mathcal{A}} = \frac{1}{2}[n_x^{\min}+n_x^{\max},\; n_y^{\min}+n_y^{\max}]^\top$ is the centroid of $\mathcal{A}$.

\subsection{Robust Optimization via Dual-Clip and Adaptive KL}
\label{subsec:dual_clip_kl}

Aggressive exploration inherently induces substantial variations in the importance-sampling ratio $r_t(\theta) = \pi_\theta(a_t \mid o_t) / \pi_{\theta_{\mathrm{old}}}(a_t \mid o_t)$, leading to oscillatory terminal overload commands $n_{y_i}$ that can severely compromise flight stability. We stabilize policy gradient optimization using a dual-clip surrogate objective combined with an adaptive KL divergence penalty.

\emph{1) Generalized Advantage Estimation:} The advantage $\hat{A}_t$ is computed via GAE with discount $\gamma_{\mathrm{RL}}$ and trace-decay $\lambda_{\mathrm{GAE}} \in [0,1]$:
\begin{align}
\hat{A}_t &= \sum_{l=0}^{T-t-1} (\gamma_{\mathrm{RL}}\lambda_{\mathrm{GAE}})^l \delta_{t+l} \label{eq:gae_A} \\
\delta_t &= r_i^{(t)} + \gamma_{\mathrm{RL}} V_\phi(s_{t+1}) - V_\phi(s_t) \label{eq:gae_delta}
\end{align}
where $\delta_t$ is the TD residual. The corresponding value target is $\hat{V}_t^{\mathrm{targ}} = \hat{A}_t + V_\phi(s_t)$.

\emph{2) Dual-Clip PPO Objective:} The standard clipped surrogate is:
\begin{equation}
\label{eq:standard_clip}
\begin{split}
L^{\mathrm{CLIP}}(\theta) = \min \Bigl( &r_t(\theta)\hat{A}_t, \\
&\operatorname{clip}\!\left(r_t(\theta), 1-\epsilon, 1+\epsilon\right)\hat{A}_t \Bigr)
\end{split}
\end{equation}
where $\epsilon \in (0,1)$ is the clipping threshold. When $\hat{A}_t < 0$, the standard clip permits $r_t(\theta) \to \infty$, causing catastrophic over-correction of overload commands. The Dual-Clip objective addresses this by introducing a floor $c > 1+\epsilon$ applied only when $\hat{A}_t < 0$:
\begin{equation}
\label{eq:dual_clip_final}
J^{\mathrm{DC}}(\theta) = \mathbb{E}_{t} \left[
\begin{cases}
\max\!\left( L^{\mathrm{CLIP}}(\theta),\; c\,\hat{A}_t \right) & \hat{A}_t < 0 \\[4pt]
L^{\mathrm{CLIP}}(\theta) & \hat{A}_t \geq 0
\end{cases}
\right]
\end{equation}
When $\hat{A}_t < 0$, the $\max$ with $c\,\hat{A}_t$ bounds $r_t(\theta) \leq c$, preventing the policy from over-penalizing actions that only marginally degrade performance. When $\hat{A}_t \geq 0$, the standard pessimistic clip suffices. The effective constraint on $r_t(\theta)$ under the dual-clip is therefore:
\begin{equation}
\label{eq:rt_bound}
r_t(\theta) \in \begin{cases} \left[\frac{1}{c},\; c\right] & \hat{A}_t < 0 \\[4pt] \left[1-\epsilon,\; 1+\epsilon\right] & \hat{A}_t \geq 0 \end{cases}
\end{equation}

\emph{3) Adaptive KL Regularization:} To continuously constrain policy deviation between updates, the KL divergence is approximated as:
\begin{equation}
\label{eq:kl_approx}
\begin{split}
\hat{D}_{\mathrm{KL}}\!\left(\pi_{\theta_{\mathrm{old}}} \| \pi_\theta\right) &= \mathbb{E}_t\!\left[\log\frac{\pi_{\theta_{\mathrm{old}}}(a_t \mid o_t)}{\pi_\theta(a_t \mid o_t)}\right] \\
&= \mathbb{E}_t\!\left[\log\pi_{\theta_{\mathrm{old}}}(a_t \mid o_t) - \log\pi_\theta(a_t \mid o_t)\right]
\end{split}
\end{equation}
The penalty coefficient $\beta_{\mathrm{KL}}$ is updated after each training epoch $k$ via:
\begin{equation}
\label{eq:kl_adapt}
\beta_{\mathrm{KL}}^{(k+1)} =
\begin{cases}
\min\!\left(1.5\,\beta_{\mathrm{KL}}^{(k)},\; \beta_{\max}\right), & \hat{D}_{\mathrm{KL}} > 2\,\delta_{\mathrm{targ}} \\[6pt]
\max\!\left(0.9\,\beta_{\mathrm{KL}}^{(k)},\; \beta_{\min}\right), & \hat{D}_{\mathrm{KL}} < 0.5\,\delta_{\mathrm{targ}} \\[6pt]
\beta_{\mathrm{KL}}^{(k)}, & \text{otherwise}
\end{cases}
\end{equation}
where $\delta_{\mathrm{targ}}$ is the target KL constraint, and $\beta_{\max}=10$, $\beta_{\min}=10^{-4}$ prevent the regularization from dominating or vanishing. When $\hat{D}_{\mathrm{KL}} > 2\delta_{\mathrm{targ}}$, the overload commands have shifted too aggressively between updates, risking structural load-factor exceedance; the $1.5\times$ increase dampens subsequent updates. When $\hat{D}_{\mathrm{KL}} < 0.5\delta_{\mathrm{targ}}$, the $0.9\times$ decay permits finer trajectory refinement.

\emph{4) Overall Loss Formulation:} The critic value loss over a mini-batch of size $B$ is:
\begin{equation}
\label{eq:value_loss}
L_V(\phi) = \frac{1}{B}\sum_{t=1}^{B} \left(V_\phi(s_t) - \hat{V}_t^{\mathrm{targ}}\right)^2
\end{equation}
Combining the dual-clip surrogate, adaptive KL penalty, value loss, and policy entropy bonus $S[\pi_\theta] = -\mathbb{E}_t[\log\pi_\theta(a_t \mid o_t)]$, the unified loss minimized by AdamW with cosine annealing is:
\begin{equation}
\label{eq:total_loss}
\mathcal{L}(\theta,\phi) = - J^{\mathrm{DC}}(\theta) + \beta_{\mathrm{KL}}\,\hat{D}_{\mathrm{KL}} + c_v\,L_V(\phi) - c_e\,S[\pi_\theta]
\end{equation}
where $c_v > 0$ and $c_e > 0$ are the value and entropy weighting coefficients. The learning rate follows a cosine annealing schedule $\eta(k) = \eta_{\min} + \frac{1}{2}(\eta_0 - \eta_{\min})(1 + \cos(\pi k / K))$, facilitating coarse trajectory optimization early in training and fine-grained terminal-guidance refinement in later epochs.

\hladd{It is worth clarifying the computational cost of the dual-clip surrogate and the adaptive-KL
penalty for latency-constrained deployment. Both mechanisms act only on the training loss
in Eq.~\eqref{eq:total_loss}: the dual-clip branch is a scalar $\max/\min$ selection on the
advantage sign, and the adaptive-KL term is a log-ratio estimate plus a scalar update of
$\beta_{\mathrm{KL}}$. Neither appears in the deployed policy. At execution each interceptor runs
only the actor forward pass, which on the deployed network ($256$ hidden units) completes in
about $0.30$~ms per decision (single interceptor) and $3.98$~ms for a $20$-agent swarm on a
server-class CPU---well inside the $50$~ms guidance cycle---while the dual-clip and adaptive-KL
operations contribute exactly $0$~ms because they are absent from this path. During offline
training the two mechanisms are evaluated once per gradient step and add only a small constant
per update ($\approx 0.013$~ms, i.e. below $0.01\%$ of a full policy update, whose cost is
dominated by the network forward/backward passes). Consequently they introduce no additional real-time overhead
on the onboard processor, and their cost during training is negligible relative to the network
forward/backward passes.}

The complete ART-MAPPO workflow, integrating all three modules described above, is illustrated in Fig.~\ref{fig:algorithm_flow}. \hladd{The corresponding training procedure under the CTDE paradigm is detailed in Algorithm~\ref{alg:artmappo}, with each line cross-referenced to the defining equation to support reproducibility.}

\begin{figure*}[!t]
\centering
\includegraphics[width=7.3in]{rlsuanfa.pdf}
\caption{Overall workflow of the ART-MAPPO algorithm for cooperative swarm interception, showing the interaction among the attention-residual network, trust-modulated exploration, and dual-clip adaptive KL optimization.}
\label{fig:algorithm_flow}
\end{figure*}

\begin{algorithm}[!t]
\color{blue}
\caption{ART-MAPPO Training for Cooperative Swarm Interception (CTDE)}
\label{alg:artmappo}
\begin{algorithmic}[1]
\REQUIRE Actor/critic params $\theta,\phi$; clip $\epsilon$, dual-clip floor $c$; target KL $\delta_{\mathrm{targ}}$; rates $\alpha_T,\rho,\tau_T$; epochs $K$
\ENSURE Trained decentralized policy $\pi_{\theta}$
\STATE Initialize $\theta,\phi$, trust $\mathcal{T}_i\!\leftarrow\!0.5$, KL weight $\beta_{\mathrm{KL}}$, running stats $\mu_R,\sigma_R^2$
\FOR{episode $k=1$ to $K$}
  \STATE {\itshape // Decentralized rollout collection}
  \FOR{each step $t$ and interceptor $i$}
    \STATE Encode $\mathbf{o}_i^{(t)}$; attention--residual features~\eqref{eq:encoder}--\eqref{eq:mha_out}, temporal state $\mathbf{c}_t$~\eqref{eq:gru}
    \STATE Form policy $\pi_{\theta_i}$~\eqref{eq:actor_dist}; sample action via trust-blended policy~\eqref{eq:exploration_policy}, \eqref{eq:guided} with $\beta(\mathcal{T}_i)=1-\mathcal{T}_i$
    \STATE Clip $a_i$ to envelope $\mathcal{A}$; step environment; store transition
  \ENDFOR
  \STATE {\itshape // Trust update from episode returns}
  \STATE Update $\mu_R,\sigma_R^2$~\eqref{eq:running_stats_mu}, \eqref{eq:running_stats_sigma}; normalize $\tilde{R}_i^{(k)}$~\eqref{eq:trust_normalized}
  \STATE Smooth trust $\mathcal{T}_i^{(k+1)}$~\eqref{eq:trust_evolution}
  \STATE {\itshape // Centralized optimization (shared critic on global state~\eqref{eq:critic_input})}
  \STATE Compute advantages $\hat{A}_t$ and value targets via GAE~\eqref{eq:gae_A}, \eqref{eq:gae_delta}
  \FOR{each mini-batch epoch}
    \STATE Evaluate dual-clip surrogate $J^{\mathrm{DC}}(\theta)$~\eqref{eq:standard_clip}, \eqref{eq:dual_clip_final}
    \STATE Estimate $\hat{D}_{\mathrm{KL}}$~\eqref{eq:kl_approx}; value loss $L_V(\phi)$~\eqref{eq:value_loss}
    \STATE Minimize total loss $\mathcal{L}(\theta,\phi)$~\eqref{eq:total_loss} with AdamW (cosine LR)
  \ENDFOR
  \STATE Adapt KL weight $\beta_{\mathrm{KL}}^{(k+1)}$~\eqref{eq:kl_adapt}
\ENDFOR
\STATE \textbf{return} decentralized policy $\pi_{\theta}$
\end{algorithmic}
\end{algorithm}

\subsection{\hladd{Theoretical Comparison with Conventional MAPPO Variants}}
\label{subsec:theory_comparison}

\hladd{Having introduced the three modules, we now position ART-MAPPO against conventional
MAPPO variants (standard MAPPO, IPPO, and single-clip PPO-penalty schemes) along the three
axes most relevant to cooperative interception: convergence guarantees, the
exploration--exploitation balance, and policy robustness. Table~\ref{tab:theory_comparison}
summarizes the comparison, and each claim is traced to the governing equation.}

\hladd{\emph{1) Convergence guarantees.} In standard MAPPO the clipped surrogate
\eqref{eq:standard_clip} is left-unbounded when $\hat{A}_t<0$: as $r_t(\theta)\!\to\!\infty$
the term $r_t(\theta)\hat{A}_t\!\to\!-\infty$, so a single mini-batch with off-policy negative
advantages can drive an arbitrarily large destructive update. In the interception setting these
are exactly the aggressive terminal-overload samples, so the standard clip admits
gradient explosions that manifest as the high inter-seed variance reported in
Section~\ref{subsec:training}. The dual-clip objective \eqref{eq:dual_clip_final}
floors the surrogate at $c\,\hat{A}_t$ and thereby confines the effective ratio to the compact
set in \eqref{eq:rt_bound}, $r_t(\theta)\in[1/c,c]$. A per-step objective that is bounded on a
compact ratio set has a Lipschitz-bounded policy gradient, so the stochastic update converges to
an $\epsilon$-stationary point of the KL-penalized surrogate under the standard stochastic-approximation
conditions for nonconvex objectives~\cite{ghadimi2013stochastic}; the bounded per-update policy shift
that this surrogate enforces is precisely the trust-region principle whose monotonic-improvement bound
motivates the design~\cite{schulman2015trust}, whereas conventional single-clip MAPPO leaves the ratio
unbounded for $\hat{A}_t\!<\!0$ and forfeits this property. We claim convergence to a stationary point
of the penalized surrogate rather than to a global optimum---the honest guarantee available to any
PPO-family method---and this bounded-gradient behavior is the mechanism behind the $2\times$ smaller
cross-seed standard deviation ($0.64$ vs.\ $1.34$) and the faster $90\%$-optimality crossing measured
empirically.}

\hladd{\emph{2) Exploration--exploitation balance.} Conventional MAPPO regulates exploration
with a \emph{fixed}, state-independent entropy bonus: every agent is pushed to explore by the
same constant amount regardless of how well it is currently performing, and the exploration is
undirected (isotropic noise on the action). ART-MAPPO instead makes the balance
\emph{performance-adaptive and structured}. The trust recursion \eqref{eq:trust_evolution} is a
contraction of modulus $(1-\alpha_T)$ toward the moving equilibrium
$\mathcal{T}_i^{\star}=\sigma(\tau_T\tilde{R}_i)$, and the exploration weight
$\beta(\mathcal{T}_i)=1-\mathcal{T}_i$ in \eqref{eq:exploration_policy} therefore shrinks for
above-average performers and grows for below-average ones---at $\tilde{R}_i=\pm2$ the allocated
exploration is $0.12$ versus $0.88$. Moreover the exploratory samples are drawn from the
tactically informative mixture \eqref{eq:guided} (proportional-navigation exploitation, envelope
probing, and residual coverage) rather than from undirected noise. The result is that
exploration is spent where it is useful and withdrawn once a competent cooperative policy is
found, which explains the smooth, monotone entropy decay observed for ART-MAPPO against the
flatter entropy profile of standard MAPPO.}

\hladd{\emph{3) Policy robustness.} Standard MAPPO controls inter-update policy drift with a
\emph{single, fixed} mechanism---either a constant clip $\epsilon$ or a constant KL weight. No
single fixed penalty is uniformly adequate: a small weight lets $\hat{D}_{\mathrm{KL}}$ spike far
outside the trust region during aggressive-exploration bursts (risking the load-factor excursions
of Section~\ref{subsec:dual_clip_kl}), while a large weight over-regularizes during calm phases
and stalls terminal-guidance refinement. The adaptive-KL controller \eqref{eq:kl_adapt} closes a
feedback loop around the target $\delta_{\mathrm{targ}}$, multiplicatively increasing
$\beta_{\mathrm{KL}}$ when $\hat{D}_{\mathrm{KL}}>2\delta_{\mathrm{targ}}$ and relaxing it when
$\hat{D}_{\mathrm{KL}}<0.5\delta_{\mathrm{targ}}$, so the realized policy shift is confined to a
bounded band across both regimes. Coupled with the two-sided ratio bound \eqref{eq:rt_bound}, this
gives ART-MAPPO a bounded, self-correcting response to the noisy, delayed observations of the
interception environment---the robustness that the Monte Carlo and hardware-in-the-loop studies
in Section~\ref{sec:simulation} confirm in closed loop. A quantitative illustration of all three
mechanisms, computed directly from Eqs.~\eqref{eq:dual_clip_final}, \eqref{eq:trust_evolution},
and \eqref{eq:kl_adapt}, is provided in the supplementary response material.}

\begin{table}[!t]
\color{blue}
\centering
\caption{\hladd{Theoretical comparison of ART-MAPPO with conventional MAPPO variants.}}
\label{tab:theory_comparison}
\renewcommand{\arraystretch}{1.25}
\resizebox{\columnwidth}{!}{%
\begin{tabular}{@{}p{2.05cm}p{2.75cm}p{2.95cm}@{}}
\toprule
\textbf{Axis} & \textbf{Conventional MAPPO} & \textbf{ART-MAPPO (this work)} \\
\midrule
Convergence guarantee &
Single clip \eqref{eq:standard_clip} unbounded for $\hat{A}_t<0$; large destructive updates possible &
Dual-clip \eqref{eq:dual_clip_final} floors objective, ratio confined to $[1/c,c]$ \eqref{eq:rt_bound}; bounded gradients \\
Exploration vs.\ exploitation &
Fixed entropy bonus; undirected, performance-independent &
Trust-modulated \eqref{eq:trust_evolution}, \eqref{eq:exploration_policy}; performance-adaptive, tactically structured \eqref{eq:guided} \\
Policy robustness &
Fixed penalty; drift not uniformly bounded across regimes &
Adaptive-KL feedback \eqref{eq:kl_adapt} confines $\hat{D}_{\mathrm{KL}}$ to a band around $\delta_{\mathrm{targ}}$ \\
Empirical effect &
Cross-seed $\sigma=1.34$; $90\%$ opt.\ at ep.\ $7654$ &
Cross-seed $\sigma=0.64$; $90\%$ opt.\ at ep.\ $5591$; $+12.3\%$ reward \\
\bottomrule
\end{tabular}%
}
\end{table}

\subsection{State Representation and Reward Design}
\label{subsec:state_reward}

\emph{1) Observation Space:} The observation vector for interceptor $i$ at time step $t$ is designed to capture the essential kinematic quantities for guidance and coordination:
\begin{equation}
\label{eq:observation_space}
\mathbf{o}_i^{(t)} = \left[o_i^{(t),1},\; o_i^{(t),2},\; o_i^{(t),3},\; o_i^{(t),4}\right]^\top \in \mathbb{R}^{n_o},\quad n_o = 4
\end{equation}
The four components are defined as:
% \begin{align}
% o_i^{(t),1} &= \begin{cases} \dfrac{V_i \dot{q}_t}{g}, & |\dot{q}_t - \dot{q}_{t-1}| \leq \Delta\dot{q}_{\max} \\[6pt] \dfrac{V_i \dot{q}_{t-1}}{g}, & |\dot{q}_t - \dot{q}_{t-1}| > \Delta\dot{q}_{\max} \end{cases} \label{eq:o1} \\[6pt]
% o_i^{(t),2} &= \gamma_{D_i} - q_i \label{eq:o2} \\[4pt]
% o_i^{(t),3} &= \left(n_{x_i}^2 + n_{y_i}^2\right) \cdot \Delta t \label{eq:o3} \\[4pt]
% o_i^{(t),4} &= \frac{r_i}{\min_{j \in \mathcal{G}_i} t_{go_j}} - V_i \label{eq:o4}
% \end{align}
\begin{align}
o_i^{(t),1} &=
\begin{cases}
\dfrac{V_i \dot{q}_i^{(t)}}{g}, &
\left|\dot{q}_i^{(t)} - \dot{q}_i^{(t-1)}\right| \leq \Delta \dot{q}_{\max}
\\[6pt]
\dfrac{V_i \dot{q}_i^{(t-1)}}{g}, &
\left|\dot{q}_i^{(t)} - \dot{q}_i^{(t-1)}\right| > \Delta \dot{q}_{\max}
\end{cases}
\label{eq:o1}
\\[6pt]
o_i^{(t),2} &= \gamma_{D_i} - q_i
\label{eq:o2}
\\[4pt]
o_i^{(t),3} &= \left(n_{x_i}^2 + n_{y_i}^2\right)\cdot \Delta t
\label{eq:o3}
\\[4pt]
o_i^{(t),4} &= \frac{r_i}{\min_{j \in \mathcal{G}_i} t_{go_j}} - V_i
\label{eq:o4}
\end{align}
where $\Delta\dot{q}_{\max} = 0.03$~rad/s is the LOS-rate noise threshold, $\mathcal{G}_i$ denotes the set of interceptors assigned to the same target as interceptor $i$, and $\Delta t$ is the simulation timestep. 
Component $o_i^{(t),1}$ is a smoothed required normal acceleration proportional to $V_i \dot{q}_i^{(t)} / g$, where the previous LOS-rate estimate is retained when the instantaneous measurement jump exceeds $\Delta \dot{q}_{\max}$.
Component $o_i^{(t),2}$ is the lead-angle deviation driving heading correction, where $q_i$ (and its derivative $\dot{q}_i$) denotes the specific line-of-sight (LOS) angle between interceptor $i$ and its dynamically assigned target. 
Component $o_i^{(t),3}$ encodes the instantaneous control energy expenditure. 
Component $o_i^{(t),4}$ is the velocity error relative to the desired approach speed $r_i / \min_j t_{go_j}$, which drives time-coordinated impact.




\emph{2) Reward Function:} The reward function incorporates five components to guide the learning process:
\begin{equation}
\label{eq:reward_function}
\begin{split}
r_i^{(t)} =& \; w_1 r_i^{\text{dist}} + w_2 r_i^{\text{angle}} + w_3 r_i^{\text{hit}} \\
&+ w_4 r_i^{\text{coord}} + w_5 r_i^{\text{energy}}
\end{split}
\end{equation}
with each component defined as:
\begin{align}
r_i^{\text{dist}} &= \exp\!\left(-\alpha_d \left\|\mathbf{p}_i^{(t)} - \mathbf{p}_{\text{tgt}}\right\|\right) \label{eq:r_dist}\\
r_i^{\text{angle}} &= -\alpha_a \left|\gamma_{D_i} - q_i\right| \label{eq:r_angle}\\
r_i^{\text{hit}} &= R_{\text{hit}} \cdot \mathbb{I}\!\left(\left\|\mathbf{p}_i^{(t)} - \mathbf{p}_{\text{tgt}}\right\| < d_{\text{kill}}\right) \label{eq:r_hit}\\
r_i^{\text{coord}} &= -\alpha_c \left| t_{go_i}^{(t)} - \bar{t}_{go}^{(t)} \right|, \quad \bar{t}_{go}^{(t)} = \frac{1}{|\mathcal{G}_i|}\sum_{j \in \mathcal{G}_i} t_{go_j}^{(t)} \label{eq:r_coord}\\
r_i^{\text{energy}} &= -\alpha_e \left(n_{x_i}^2 + n_{y_i}^2\right) \label{eq:r_energy}
\end{align}
where $\mathbf{p}_i^{(t)} = [x_i, y_i]^\top$ is the interceptor position, $\mathbf{p}_{\text{tgt}}$ is the target position, $d_{\text{kill}}$ is the lethal radius, $R_{\text{hit}} > 0$ is the terminal success reward, and $w_k$, $\alpha_k$ are weighting coefficients. The distance reward $r_i^{\text{dist}} \in (0,1]$ provides dense, bounded feedback for target approach. The coordination reward $r_i^{\text{coord}}$ penalizes deviation from the group mean time-to-go $\bar{t}_{go}^{(t)}$, encouraging simultaneous impact; this term is computed by the centralized critic during training and obtained via inter-agent communication during execution. The energy term $r_i^{\text{energy}}$ regularizes control effort to prevent actuator saturation.

\hladd{To characterize how the reward weights shape cooperative behavior, we perform a one-at-a-time sensitivity analysis in which each $w_k$ is scaled over $[0.1,2.0]\times$ its nominal value while the others are held fixed, tracking the four evaluation metrics of Section~\ref{subsec:eval}. The results confirm that each term primarily governs its intended tactical dimension. The distance weight $w_1$ is the dominant driver of terminal miss distance (normalized local sensitivity $|\partial \ln E_{miss}/\partial \ln w_1| \approx 0.76$), improving $E_{miss}$ and interception success rate at the expense of a mild rise in normal overload. The energy weight $w_5$ most strongly regulates terminal overload $E_n$ (sensitivity $\approx 0.71$): increasing it suppresses maneuver aggressiveness but degrades miss distance and success rate if overweighted. The coordination weight $w_4$ dominates temporal coordination $E_{co\text{-}time}$ (sensitivity $\approx 0.78$), tightening simultaneous impact at a slight precision cost, while the hit weight $w_3$ and angular weight $w_2$ mainly reinforce success rate and heading alignment. Because the metric--weight coupling is monotone and well separated, the reward weights can be tuned to shift the operating point toward precision, safety, or temporal-coordination priorities without destabilizing the others; the nominal setting adopted in this work balances these objectives.}

\subsection{Evaluation Criteria}
\label{subsec:eval}

To quantitatively evaluate the proposed cooperative interception method, we adopt four metrics that capture temporal coordination, terminal load safety, guidance accuracy, and operational efficiency. The first three metrics are formally defined as follows.
\begin{equation}
\label{E1}
E_{co\text{-}time} = \frac{1}{m}\sum_{i=1}^{m} \left| t_{go_i} - \frac{1}{m}\sum_{j=1}^{m} t_{go_j} \right|
\end{equation}
\begin{equation}
\label{E2}
E_{n} = \frac{1}{m}\sum_{i=1}^{m} |n_{y_i}|
\end{equation}
\begin{equation}
\label{E3}
E_{miss} = \frac{1}{m}\sum_{i=1}^{m} d_{tf_i}
\end{equation}

Here, Eq.~\eqref{E1} measures the mean deviation of individual time-to-go values from the group mean, thereby assessing temporal coordination among interceptors assigned to the same target. Eq.~\eqref{E2} gives the average terminal normal overload, which reflects control effort and maneuver aggressiveness during the terminal guidance phase. Eq.~\eqref{E3} denotes the mean terminal miss distance, where $d_{tf_i}$ is the relative distance between interceptor $i$ and its assigned target at the terminal engagement time $t_{f_i}$, serving as a direct indicator of guidance accuracy. The fourth metric, coordinated interception duration, quantifies operational efficiency and is reported as the mean elapsed time between the earliest and latest terminal engagements within each coordinated group. All four metrics are reported in the simulation results to provide a comprehensive assessment of effectiveness, safety, precision, and efficiency.










%% ===========================================================
\section{Simulation Verification}
\label{sec:simulation}

\subsection{Experimental Setup and Implementation Details}
\label{subsec:setup}

All training and evaluation were performed on an Intel Core i7-9700F CPU / Nvidia RTX 3090
GPU server running Ubuntu 18.04.
The simulation step size was 0.05\,s.
Table~\ref{table2} summarizes the key hyperparameters of ART-MAPPO.

\begin{table}[!t]
\centering
\renewcommand{\arraystretch}{1.4}
\caption{ART-MAPPO Algorithm Parameters}
\label{table2}
\begin{tabular}{ll}
\hline\hline
Parameter & Value \\ \hline
PPO clipping factor ($\epsilon$) & 0.2 \\
Dual-clip parameter ($c$) & 3.0 \\
Target KL divergence ($\delta_{\text{targ}}$) & 0.02 \\
Entropy bonus coefficient ($c_e$) & 0.01 \\
GAE parameter ($\lambda_{\mathrm{GAE}}$) & 0.97 \\
\hladd{Discount factor ($\gamma_{\mathrm{RL}}$)} & \hladd{0.99} \\
LR schedule & Cosine Annealing \\
Optimizer & AdamW \\
Mini-batch number & 4 \\
Buffer size & 4096 \\
Learning rate & $3\times10^{-4}$ \\
Attention heads ($H$) & 4 \\
Residual blocks ($L$) & 2 \\ \hline
\end{tabular}
\end{table}

\hladd{The behavior of the learned policy is governed jointly by the reward-shaping weights
$w_1$--$w_5$ in \eqref{eq:reward_function} and by the training hyperparameters in
Table~\ref{table2}. The weights act as tactical dials: a one-at-a-time sweep over
$[0.1,2]\times$ their nominal values shifts each evaluation metric along the physical
trade-off it controls (distance weight $w_1\!\to\!$ terminal miss, coordination weight
$w_4\!\to\!$ time-to-go dispersion, energy weight $w_5\!\to\!$ terminal overload) while the
interception success rate stays above $95\%$ throughout, so no single weight is a knife-edge
on effectiveness and $(w_1,w_4)\!\approx\!(1.0,0.6)$ is a balanced compromise. The training
hyperparameters, by contrast, affect the outcome only weakly: sweeping the learning rate
($0.1$--$10\times$), clip factor ($0.05$--$0.40$), KL target ($0.1$--$10\times$), discount,
GAE parameter, and entropy coefficient each keeps the converged success rate within $1\%$ of
its nominal value over a broad band, with the cross-seed reward variance rising only outside
that band. This flatness is by design: the dual-clip surrogate \eqref{eq:dual_clip_final}
bounds the policy ratio irrespective of $\epsilon$, and the adaptive-KL controller
\eqref{eq:kl_adapt} drives the realized divergence back to $\delta_{\text{targ}}$ regardless of
the raw step size, so learning quality is set by the problem-specific reward geometry rather
than by fine hyperparameter tuning.}


The simulation environment is depicted in Fig.~\ref{chushi}.
$A_{\text{UAV}}$ are randomly spawned within an annular region of radii $D_1$--$D_2$\,m with speed
$V_A\in[10,40]$\,m/s. To simulate a directional swarm attack toward the high-value target area, their initial heading angles $\gamma_A$ are uniformly distributed within $[q_A-30^\circ, q_A+30^\circ]$, where $q_A$ explicitly denotes the initial line-of-sight (LOS) angle from the specific attacker to the center of the defended area.
$D_{\text{UAV}}$ are uniformly distributed within a disk of radius $D_3$\,m with
$V_D\in[10,40]$\,m/s and $\gamma_D\in[0,2\pi]$.
\hladd{In the three-dimensional scenario, the defenders start near the defended center at
low altitude and the attackers start at higher altitude, so that every engagement requires
simultaneous horizontal and vertical closure; the initial flight-path angles $\theta$ are set
from the initial defender--target elevation and the attackers descend toward the defended
region while executing the case-specific three-dimensional maneuvers described below.}


\begin{figure}[!t]
\centering
\includegraphics[width=3.0in]{chushi01.pdf}
\caption{Simulation environment configuration for training.}
\label{chushi}
\end{figure}

A representative scenario of 20 $D_{\text{UAV}}$ versus 8 $A_{\text{UAV}}$ is considered, with $D_1=\hldel{100}\hladd{1500}$\,m,
$D_2=\hldel{1500}\hladd{1600}$\,m, and $D_3=\hldel{1600}\hladd{100}$\,m\hladd{, so that the attackers
spawn in a far annulus ($D_1$--$D_2$) and descend toward the defended center while the defenders
start clustered within a radius-$D_3$ disk around it, consistent with Fig.~\ref{chushi}}.
Table~\ref{chushicanshu} summarizes the UAV kinematic parameters.

\begin{table*}[!t]
\centering
\caption{UAV Parameters for Attackers and Defenders}
\label{chushicanshu}
\renewcommand{\arraystretch}{1.6}
\begin{tabular}{ccccc}
\hline\hline
 & Initial position \hldel{$(x,y)$}\hladd{$(x,y,z)$} & Initial velocity $V$ & Initial heading $\gamma$ & Maneuverability \hldel{$(n_x,n_y,V)$}\hladd{$(n_x,n_y,n_z)$} \\ \hline
$D_{\text{UAV}}$ & $r\sim U[0,D_3]$, $\theta\sim U[0,2\pi]$\hladd{, $z\!\approx\!0$} & $V_D\sim U[10,40]$\,m/s & $\gamma_D\sim U[0,2\pi]$ & $n_x\in[-0.1g,1g]$, $n_y\in[-1g,1g]$\hladd{, $n_z\in[-1g,1g]$} \\
$A_{\text{UAV}}$ & $r\sim U[D_1,D_2]$, $\theta\sim U[0,2\pi]$\hladd{, elevated $z$} & $V_A\sim U[10,40]$\,m/s & $\gamma_A\sim U[q_A-30^\circ,q_A+30^\circ]$ & $n_x=0$, $n_y\in[-1g,1g]$\hladd{, $n_z\in[-1g,1g]$} \\ \hline
\end{tabular}
\end{table*}

Two dynamic engagement scenarios are designed to evaluate cooperative guidance performance.
\hladd{A third scenario, Case~3, is reserved as a held-out generalization test: its geometry and
attack pattern are never seen during training, and it is used in
Section~\ref{subsec:generalization} to assess zero-shot transfer of the trained policy.}
The scenario configurations are listed in Table~\ref{changjingqingkuang}.

\begin{table}[!t]
\centering
\renewcommand{\arraystretch}{1.4}
\caption{Combat Scenario Settings}
\label{changjingqingkuang}
\begin{tabular}{cccc}
\hline\hline
 & $N_{D\text{-UAV}}$ & $N_{A\text{-UAV}}$ & $A_{\text{UAV}}$ Strategy \\ \hline
Case 1 & 20 & 8 & Hybrid Evasion \\
Case 2 & 20 & 8 & Sinusoidal Maneuver \\
\hladd{Case 3} & \hladd{20} & \hladd{8} & \hladd{Unseen Hybrid Maneuver} \\ \hline
\end{tabular}
\end{table}

%% ==== Re2_4 / Re3_7 / Re3_10: environmental modeling, uncertainty, and HIL ====
\hladd{\emph{1) Communication Delay and Sensing Uncertainty:} Beyond the UAV dynamics, the
environment explicitly models delayed communication and imperfect sensing. Every defender
receives the relative position and velocity packets used to build its observation
Eq.~\eqref{eq:observation_space} only after a finite transport delay, and each Cartesian
component is corrupted by zero-mean Gaussian error,}
\begin{equation}
\label{eq:noise_model}
{\color{blue}
\widehat{\mathbf{p}}(t)=\mathbf{p}(t-\tau_s)+\boldsymbol{\eta}_p,\quad
\widehat{\mathbf{v}}(t)=\mathbf{v}(t-\tau_s)+\boldsymbol{\eta}_v,}
\end{equation}
\hladd{with $\boldsymbol{\eta}_p\sim\mathcal{N}(\mathbf{0},\sigma_p^2\mathbf{I}_3)$ and
$\boldsymbol{\eta}_v\sim\mathcal{N}(\mathbf{0},\sigma_v^2\mathbf{I}_3)$. Range, altitude
difference, LOS direction, closing speed, heading, pitch, velocity-to-LOS angle, and
time-to-go are all recomputed from these delayed noisy Cartesian states, so no independent
noise is added to the derived quantities. Command packets are additionally subject to a
first-order actuation response and a small probability of loss, in which case the last
successfully received command is held. Table~\ref{tab:uncertainty_settings} lists the settings
in physical units.}

\begin{table}[!t]
\centering
\caption{\hladd{Communication Delay and Observation-Uncertainty Settings}}
\label{tab:uncertainty_settings}
\renewcommand{\arraystretch}{1.2}
\footnotesize
\begin{tabular}{@{}p{0.30\linewidth}p{0.22\linewidth}p{0.38\linewidth}@{}}
\hline\hline
{\color{blue}Observed quantity} & {\color{blue}Delay / response} & {\color{blue}Error distribution} \\
\hline
{\color{blue}Position $x,y,z$} & {\color{blue}$\tau_s=100$\,ms} & {\color{blue}Mean $0$; per-component $\sigma_p=3$\,m} \\
{\color{blue}Velocity $v_x,v_y,v_z$} & {\color{blue}$\tau_s=100$\,ms} & {\color{blue}Mean $0$; per-component $\sigma_v=0.3$\,m/s} \\
{\color{blue}Derived geometry, $t_{go}$} & {\color{blue}$100$\,ms} & {\color{blue}Recomputed from $\widehat{\mathbf{p}},\widehat{\mathbf{v}}$} \\
{\color{blue}Command channel} & {\color{blue}$50$\,ms} & {\color{blue}First-order hold; $1\%$ packet loss} \\
\hline\hline
\end{tabular}
\end{table}

\hladd{\emph{2) Hardware-in-the-Loop Semi-Physical Validation:} To move beyond a pure software
study and improve practical credibility, the three-dimensional experiment is executed as a
hardware-in-the-loop (HIL) semi-physical setup. Five interceptor policy endpoints ($D_0$--$D_4$)
run on separate NVIDIA Jetson NX embedded computers, each maintaining its own recurrent state,
observation-delay line, and actuator-delay line, exactly as they would on a physical
interceptor. The remaining interceptors and all attackers evolve in the mathematical simulator,
which exchanges state and command packets with the NX nodes over the network under the delay,
noise, and packet-loss channel of Table~\ref{tab:uncertainty_settings}. This closed loop
therefore validates the decentralized-execution assumption of the CTDE design on real
processors and real inter-node communication, rather than in a single monolithic simulation.
The platform architecture is shown in Fig.~\ref{fig:hil_platform}: a server (Ubuntu~22.04,
Core~i9-10900K, RTX~3080) runs the engagement-state update, the attacker cluster, and the
software policies of the remaining interceptors ($D_5$--$D_n$), while five Jetson NX nodes
(one per hardware policy endpoint $D_0$--$D_4$) execute online inference and their target-assignment
modules; server and embedded sides are interconnected through a physical TCP Ethernet link over a
TP-LINK switch, carrying the observation downlink and the action/target-assignment uplink.}

\begin{figure}[!t]
\centering
\includegraphics[width=0.45\textwidth]{tase_HIL.png}
\caption{\hladd{Hardware-in-the-loop, semi-physical validation platform. The server side
(Ubuntu~22.04, Core~i9-10900K, RTX~3080) hosts the engagement-state update, the attacker
cluster, and the software policies of interceptors $D_5$--$D_n$; the five hardware policy
endpoints $D_0$--$D_4$ run online inference and target assignment on separate NVIDIA Jetson~NX
embedded computers. The two sides communicate over a physical TCP Ethernet link (TP-LINK switch),
exchanging the observation downlink and the action/target-assignment uplink under the delay,
noise, and packet-loss channel of Table~\ref{tab:uncertainty_settings}. In the diagram the agent
index $A_k$ labels the policy of interceptor $D_k$, i.e. $A_0$--$A_4$ are the hardware endpoints
$D_0$--$D_4$ and $A_5$--$A_n$ are the software-policy interceptors $D_5$--$D_n$.}}
\label{fig:hil_platform}
\end{figure}





\subsection{Target Assignment Performance and Spatial Topology}
\label{subsec:assignment_results}

To evaluate the IDBO algorithm, it is benchmarked against PSO, GWO, SSA, BOA, and standard DBO. 

Fig.~\ref{fig:convergence} illustrates the statistical performance in terms of average cost convergence over iterations. IDBO significantly outperforms the baselines, reaching the lowest steady-state cost within 100~iterations. \begin{figure}[!ht] 
\centering 
\includegraphics[width=3.25in]{convergence.png} 
\caption{Average cost convergence over iterations for target assignment algorithms.} 
\label{fig:convergence} 
\end{figure}

Furthermore, the violin plot in Fig.~\ref{fig:violin} evaluates the algorithmic robustness across multiple independent runs. It shows IDBO's compact distribution at the lowest median cost, highlighting its superior stability and consistency compared to the other heuristic methods.

\begin{figure}[!ht] 
\centering 
\includegraphics[width=3.2in]{violin.png} 
\caption{Violin plot of total cost distribution evaluating algorithm robustness.} 
\label{fig:violin} 
\end{figure}

The optimal assignment mapping generated by the algorithm is visualized in Fig.~\ref{fig:assignment_topology}. Attackers $A_1$--$A_4$ are each allocated three defenders, while $A_5$--$A_8$ are intercepted by two defenders each, perfectly respecting the per-target capacity constraint $L_{\max}$ and ensuring sufficient numerical superiority for cooperative strikes.

\begin{figure}[!ht]
\centering
\includegraphics[width=3.2in]{target_assignment_topology.png}
\caption{Visualized spatial topology of the target assignment results.}
\label{fig:assignment_topology}
\end{figure}








\subsection{RL Model Training and Convergence Analysis}
\label{subsec:training}

To evaluate the learning efficiency and robustness of the proposed ART-MAPPO, comparative training experiments were conducted against standard MAPPO, IPPO, IA2C, and IQL. All models were trained over 10,000 episodes with a maximum horizon of \hldel{25}\hladd{2500} steps per episode. To ensure statistical significance and eliminate the bias of favorable network initialization, all learning curves in Fig.~\ref{fig:training_convergence} are averaged across 5 independent random seeds, with the shaded regions denoting the standard deviation.

As illustrated in Fig.~\ref{fig:training_convergence}(a), ART-MAPPO demonstrates a superior learning paradigm, achieving the highest asymptotic normalized reward of $8.97 \pm 0.64$ (averaged over the final 500 episodes). This represents a notable performance improvement of 8.2\%, 12.3\%, 14.4\%, and 29.8\% over IQL, MAPPO, IPPO, and IA2C, respectively. The advantage of ART-MAPPO is particularly pronounced during the intermediate training phase. By episode 4,000, ART-MAPPO establishes a significant performance leap, outperforming standard MAPPO by 43.6\% and IQL by 16.6\%, highlighting its ability to rapidly escape local optima early in the training process.

Furthermore, the proposed method exhibits accelerated convergence alongside enhanced training stability. When evaluated on a 100-episode smoothed curve, ART-MAPPO reaches 90\% of its optimal asymptotic performance by episode 5,591. In contrast, the most competitive baseline, IQL, requires 6,477 episodes (a 13.7\% delay), while standard MAPPO lags significantly, requiring 7,654 episodes. Regarding algorithmic stability, ART-MAPPO maintains a tight variance across random seeds ($\sigma = 0.64$), whereas standard MAPPO exhibits a standard deviation of 1.34 (approximately twice as large). This confirms that the proposed enhancements effectively mitigate the high-variance gradient updates typical in multi-agent policy-based methods.

These macroscopic performance gains are directly supported by the network-level metrics. The critic loss (Fig.~\ref{fig:training_convergence}(b)) of ART-MAPPO converges to the lowest steady-state error with minimal oscillation, a direct consequence of the dual-clip mechanism and adaptive KL penalty. Concurrently, the policy entropy (Fig.~\ref{fig:training_convergence}(c)) demonstrates an optimal, gradual decay trajectory, validating that the trust-modulated exploration mechanism successfully balances early-stage extensive state-space exploration with late-stage strict policy exploitation.

\begin{figure}[!t]
\centering
\subfloat[Training Reward]{\includegraphics[width=3.0in]{fig_a_reward.png}}
\hfill
\subfloat[Critic Loss]{\includegraphics[width=3.0in]{fig_b_critic_loss.png}}
\hfill
\subfloat[Policy Entropy]{\includegraphics[width=3.0in]{fig_c_entropy.png}}
\caption{Comparative training curves of ART-MAPPO and baseline algorithms over 10,000 episodes. Solid lines and shaded regions represent the mean and standard deviation across 5 independent random seeds, respectively. Notably, ART-MAPPO achieves the highest asymptotic reward (+12.3\% over standard MAPPO) and the fastest convergence rate (reaching 90\% optimality by episode 5,591), while demonstrating significantly enhanced training stability.}
\label{fig:training_convergence}
\end{figure}


%%%%-------------------------------------------------------------------------------------------

\subsection{Cooperative Interception against Maneuvering Targets}
\label{subsec:cases}

Based on the optimal assignment topology, we evaluate \hldel{the terminal}\hladd{three-dimensional}
cooperative guidance \hldel{performance against}\hladd{in} two maneuvering scenarios\hldel{, each benchmarked against classical PN}\hladd{. ART-MAPPO and
the capacity-matched PN reference share the same initial states, target assignment, hit radius,
speed limits, and load-factor envelope. For each method and case we report the
three-dimensional trajectory, its horizontal projection, all three load commands with their
attitude/speed responses, the time-to-go histories, and the terminal arrival-time differences}.







\subsubsection{Case~1: Interception against Evasive Maneuvering Targets}

In \hldel{this scenario, the attacking UAVs ($A_{\text{UAV}}$)}\hladd{Case~1, the attackers}
initially execute a \hldel{constant-$g$}\hladd{constant-load} evasive maneuver and transition to
\hldel{proportional navigation (PN)}\hladd{PN} guidance after 15\,s. \hldel{The defending swarm is tasked with neutralizing these highly dynamic threats utilizing the pre-computed assignment topology.}\hladd{This
creates simultaneous horizontal and vertical closing demands while the defenders must preserve
the assignment groups and coordinate their terminal arrival times.}

\hladd{\emph{1) ART-MAPPO:}} \hldel{The following figure illustrates the 2D engagement trajectories generated by the proposed ART-MAPPO algorithm. Guided by the intelligent policy, the defending swarm effectively divides into eight coordinated combat units. Despite the evasive maneuvers of the targets, ART-MAPPO generates remarkably smooth and direct pursuit curves, enabling all interceptors to successfully converge on their respective targets without unnecessary detours.}%
\hladd{Fig.~\ref{fig:case1_mappo_traj_3d} shows the complete spatial engagement: starting below the
attackers, the defenders climb from the defended region and close on the eight elevated target
trajectories. The horizontal projection in Fig.~\ref{fig:case1_mappo_traj_xy} preserves the
original plan-view interpretation and confirms that all 20 defenders converge on the targets
specified by the assignment topology.}

\begin{figure}[!ht]
\centering
\includegraphics[width=3.15in]{\vfig{mappo/mappo_nopn_trajectory_3d.png}}
\caption{\hladd{Three-dimensional engagement trajectories of the defending swarm using ART-MAPPO
in Case~1.}}
\label{fig:case1_mappo_traj_3d}
\end{figure}

\begin{figure}[!ht]
\centering
\includegraphics[width=3.15in]{\vfig{mappo/mappo_nopn_trajectory.png}}
\caption{\hladd{Horizontal projection of the ART-MAPPO engagement trajectories in Case~1.}}
\label{fig:case1_mappo_traj_xy}
\end{figure}

\hladd{The command and state histories are grouped into three physically related pairs in
Figs.~\ref{fig:case1_mappo_yaw}--\ref{fig:case1_mappo_axial}. The horizontal normal command $n_y$
and heading $\gamma_D$ form the yaw-plane response, while the vertical command $n_z$ and pitch
$\theta_D$ capture the vertical response that is absent from a planar model. ART-MAPPO keeps
$n_y$ inside the $\pm1$\,g envelope and uses a substantially smaller vertical command for most
of the engagement, and the smooth pitch histories confirm gradual altitude capture. The axial
command $n_x$ and speed $V_D$ show how the policy regulates closing speed without decoupling
temporal coordination from the two normal-load channels.}

\begin{figure}[!ht]
\centering
\includegraphics[width=3.1in]{\vfig{mappo/standalone_duav_legend.png}}\\[-2pt]
\subfloat[Normal load $n_y$]{\includegraphics[width=3.15in]{\vfig{mappo/mappo_nopn_ny.png}}}\\[-2pt]
\subfloat[Heading $\gamma_D$]{\includegraphics[width=3.15in]{\vfig{mappo/mappo_nopn_heading.png}}}
\caption{\hladd{Yaw-plane command and response of ART-MAPPO in Case~1.}}
\label{fig:case1_mappo_yaw}
\end{figure}

\begin{figure}[!ht]
\centering
\includegraphics[width=3.1in]{\vfig{mappo/standalone_duav_legend.png}}\\[-2pt]
\subfloat[Vertical load $n_z$]{\includegraphics[width=3.15in]{\vfig{mappo/mappo_nopn_nz.png}}}\\[-2pt]
\subfloat[Pitch $\theta_D$]{\includegraphics[width=3.15in]{\vfig{mappo/mappo_nopn_pitch.png}}}
\caption{\hladd{Pitch-plane command and response of ART-MAPPO in Case~1.}}
\label{fig:case1_mappo_pitch}
\end{figure}

\begin{figure}[!ht]
\centering
\includegraphics[width=3.1in]{\vfig{mappo/standalone_duav_legend.png}}\\[-2pt]
\subfloat[Axial load $n_x$]{\includegraphics[width=3.15in]{\vfig{mappo/mappo_nopn_nx.png}}}\\[-2pt]
\subfloat[Velocity $V_D$]{\includegraphics[width=3.15in]{\vfig{mappo/mappo_nopn_velocity.png}}}
\caption{\hladd{Axial command and speed response of ART-MAPPO in Case~1.}}
\label{fig:case1_mappo_axial}
\end{figure}

\hladd{The $t_{go}$ histories in Fig.~\ref{fig:case1_mappo_tgo} converge within the eight target
groups despite the three-dimensional engagement geometry, and the terminal summary in
Fig.~\ref{fig:case1_mappo_sync} reports a maximum within-group arrival spread of 0.25\,s, well
below the 0.5-s synchronization criterion.}

\begin{figure}[!ht]
\centering
\includegraphics[width=3.2in]{\vfig{mappo/mappo_nopn_tgo.png}}
\caption{\hldel{Time-to-go ($t_{go}$) profiles of the defending swarm using ART-MAPPO in Case~1. Interceptors assigned to the same target actively synchronize their remaining flight times to guarantee a simultaneous saturation strike.}\hladd{Time-to-go ($t_{go}$) profiles of the defending swarm using ART-MAPPO in Case~1.}}
\label{fig:case1_mappo_tgo}
\end{figure}

\begin{figure}[!ht]
\centering
\includegraphics[width=3.0in]{\vfig{mappo/mappo_nopn_time_sync.png}}
\caption{\hldel{Terminal arrival time differences $\Delta t$ within each coordinated interceptor group under ART-MAPPO. The near-zero $\Delta t$ values explicitly validate the successful execution of temporally synchronized saturation strikes.}\hladd{Terminal arrival time differences $\Delta t$ within each ART-MAPPO interceptor group in Case~1.}}
\label{fig:case1_mappo_sync}
\end{figure}

\hladd{\emph{2) PN baseline:}} \hladd{The capacity-matched PN reference is shown in
Figs.~\ref{fig:case1_pn_traj_3d} and~\ref{fig:case1_pn_traj_xy}. PN also completes the selected
engagement, but its purely reactive commands produce visibly different spatial curvature and
control transients. Identical command limits are imposed on both methods, so the comparison does
not attribute physically unavailable load authority to the baseline. The corresponding
yaw-plane, pitch-plane, and axial responses in Figs.~\ref{fig:case1_pn_yaw}--\ref{fig:case1_pn_axial}
show more direct use of the available bounds, while the pitch response and speed histories remain
within the common vehicle envelope, providing a physically matched three-dimensional reference.
Its $t_{go}$ histories (Fig.~\ref{fig:case1_pn_tgo}) also converge, with a maximum terminal
group spread of 0.40\,s (Fig.~\ref{fig:case1_pn_sync}) versus 0.25\,s for ART-MAPPO; both satisfy
the 0.5-s criterion while the learned policy gives the tighter representative synchronization.}

\begin{figure}[!ht]
\centering
\includegraphics[width=3.15in]{\vfig{pn/pn_nopn_trajectory_3d.png}}
\caption{\hladd{Three-dimensional engagement trajectories of the defending swarm using the PN
baseline in Case~1.}}
\label{fig:case1_pn_traj_3d}
\end{figure}

\begin{figure}[!ht]
\centering
\includegraphics[width=3.15in]{\vfig{pn/pn_nopn_trajectory.png}}
\caption{\hladd{Horizontal projection of the PN engagement trajectories in Case~1.}}
\label{fig:case1_pn_traj_xy}
\end{figure}

\begin{figure}[!ht]
\centering
\includegraphics[width=3.1in]{\vfig{pn/standalone_duav_legend.png}}\\[-2pt]
\subfloat[Normal load $n_y$]{\includegraphics[width=3.15in]{\vfig{pn/pn_nopn_ny.png}}}\\[-2pt]
\subfloat[Heading $\gamma_D$]{\includegraphics[width=3.15in]{\vfig{pn/pn_nopn_heading.png}}}
\caption{\hladd{Yaw-plane command and response of the capacity-matched PN baseline in Case~1.}}
\label{fig:case1_pn_yaw}
\end{figure}

\begin{figure}[!ht]
\centering
\includegraphics[width=3.1in]{\vfig{pn/standalone_duav_legend.png}}\\[-2pt]
\subfloat[Vertical load $n_z$]{\includegraphics[width=3.15in]{\vfig{pn/pn_nopn_nz.png}}}\\[-2pt]
\subfloat[Pitch $\theta_D$]{\includegraphics[width=3.15in]{\vfig{pn/pn_nopn_pitch.png}}}
\caption{\hladd{Pitch-plane command and response of the capacity-matched PN baseline in Case~1.}}
\label{fig:case1_pn_pitch}
\end{figure}

\begin{figure}[!ht]
\centering
\includegraphics[width=3.1in]{\vfig{pn/standalone_duav_legend.png}}\\[-2pt]
\subfloat[Axial load $n_x$]{\includegraphics[width=3.15in]{\vfig{pn/pn_nopn_nx.png}}}\\[-2pt]
\subfloat[Velocity $V_D$]{\includegraphics[width=3.15in]{\vfig{pn/pn_nopn_velocity.png}}}
\caption{\hladd{Axial command and speed response of the capacity-matched PN baseline in Case~1.}}
\label{fig:case1_pn_axial}
\end{figure}

\begin{figure}[!ht]
\centering
\includegraphics[width=3.2in]{\vfig{pn/pn_nopn_tgo.png}}
\caption{\hladd{Time-to-go ($t_{go}$) profiles of the defending swarm using the PN baseline in
Case~1.}}
\label{fig:case1_pn_tgo}
\end{figure}

\begin{figure}[!ht]
\centering
\includegraphics[width=3.0in]{\vfig{pn/pn_nopn_time_sync.png}}
\caption{\hladd{Terminal arrival time differences $\Delta t$ within each PN interceptor group in
Case~1.}}
\label{fig:case1_pn_sync}
\end{figure}







\subsubsection{Case~2: Interception against Continuous Weaving Targets}

\hldel{This scenario evaluates the robustness of the proposed framework against the}\hladd{Case~2 uses a}
sinusoidal penetration maneuver~\cite{Zarchan1995969}\hldel{, which represents a highly agile and continuous weaving threat. Such maneuvers induce high-frequency oscillations in the line-of-sight (LOS) rate, severely challenging both the guidance stability and the temporal coordination of the defending swarm.}\hladd{. Its continuous horizontal and
vertical LOS variation provides a more persistent disturbance than the Case~1 transition and
tests whether group timing can be maintained while all three command channels respond.}

\hladd{\emph{1) ART-MAPPO:}} \hldel{The 2D engagement trajectories of the ART-MAPPO defending swarm are illustrated in the following figure. Despite the continuous lateral oscillations of the $A_{\text{UAV}}$, ART-MAPPO explicitly predicts and accommodates the target's weaving motion. The defending swarm maintains all eight coordinated combat units and successfully intercepts every target with smooth pursuit curves.}%
\hladd{The three-dimensional paths in Fig.~\ref{fig:case2_mappo_traj_3d} show that all eight groups
maintain spatial closure against the weaving targets, and the horizontal projection in
Fig.~\ref{fig:case2_mappo_traj_xy} exposes the lateral curvature without hiding the altitude
evolution shown in the 3-D view.}

\begin{figure}[!ht]
\centering
\includegraphics[width=3.15in]{\vfig{mappo/mappo_sin_trajectory_3d.png}}
\caption{\hladd{Three-dimensional ART-MAPPO engagement trajectories against continuous weaving
targets in Case~2.}}
\label{fig:case2_mappo_traj_3d}
\end{figure}

\begin{figure}[!ht]
\centering
\includegraphics[width=3.15in]{\vfig{mappo/mappo_sin_trajectory.png}}
\caption{\hladd{Horizontal projection of the ART-MAPPO engagement trajectories in Case~2.}}
\label{fig:case2_mappo_traj_xy}
\end{figure}

\hladd{As shown in Figs.~\ref{fig:case2_mappo_yaw}--\ref{fig:case2_mappo_axial}, ART-MAPPO adapts
$n_y$ to the continuous lateral disturbance while using $n_z$ to regulate vertical closure. The
heading and pitch histories remain smooth relative to the command variations, and all load
channels remain inside the prescribed envelope. The axial command raises or preserves closing
speed without eliminating the control authority required in the yaw and pitch planes.}

\begin{figure}[!ht]
\centering
\includegraphics[width=3.1in]{\vfig{mappo/standalone_duav_legend.png}}\\[-2pt]
\subfloat[Normal load $n_y$]{\includegraphics[width=3.15in]{\vfig{mappo/mappo_sin_ny.png}}}\\[-2pt]
\subfloat[Heading $\gamma_D$]{\includegraphics[width=3.15in]{\vfig{mappo/mappo_sin_heading.png}}}
\caption{\hladd{Yaw-plane command and response of ART-MAPPO in Case~2.}}
\label{fig:case2_mappo_yaw}
\end{figure}

\begin{figure}[!ht]
\centering
\includegraphics[width=3.1in]{\vfig{mappo/standalone_duav_legend.png}}\\[-2pt]
\subfloat[Vertical load $n_z$]{\includegraphics[width=3.15in]{\vfig{mappo/mappo_sin_nz.png}}}\\[-2pt]
\subfloat[Pitch $\theta_D$]{\includegraphics[width=3.15in]{\vfig{mappo/mappo_sin_pitch.png}}}
\caption{\hladd{Pitch-plane command and response of ART-MAPPO in Case~2.}}
\label{fig:case2_mappo_pitch}
\end{figure}

\begin{figure}[!ht]
\centering
\includegraphics[width=3.1in]{\vfig{mappo/standalone_duav_legend.png}}\\[-2pt]
\subfloat[Axial load $n_x$]{\includegraphics[width=3.15in]{\vfig{mappo/mappo_sin_nx.png}}}\\[-2pt]
\subfloat[Velocity $V_D$]{\includegraphics[width=3.15in]{\vfig{mappo/mappo_sin_velocity.png}}}
\caption{\hladd{Axial command and speed response of ART-MAPPO in Case~2.}}
\label{fig:case2_mappo_axial}
\end{figure}

\hladd{Despite the continuous LOS disturbance, the ART-MAPPO $t_{go}$ curves converge within their
assigned groups (Fig.~\ref{fig:case2_mappo_tgo}). The largest terminal arrival-time difference is
0.30\,s in Fig.~\ref{fig:case2_mappo_sync}, demonstrating that the timing objective remains
active in the three-dimensional engagement.}

\begin{figure}[!ht]
\centering
\includegraphics[width=3.2in]{\vfig{mappo/mappo_sin_tgo.png}}
\caption{\hldel{Time-to-go ($t_{go}$) profiles of the defending swarm using ART-MAPPO in Case~2. Interceptors actively synchronize their remaining flight times despite the continuous high-G maneuvers of the targets.}\hladd{Time-to-go ($t_{go}$) profiles of the defending swarm using ART-MAPPO in Case~2.}}
\label{fig:case2_mappo_tgo}
\end{figure}

\begin{figure}[!ht]
\centering
\includegraphics[width=3.0in]{\vfig{mappo/mappo_sin_time_sync.png}}
\caption{\hldel{Terminal arrival time differences $\Delta t$ for Case~2. The remarkably low variance ($\le 0.30$\,s) proves that ART-MAPPO preserves strict temporal synchronization even against continuous weaving targets.}\hladd{Terminal arrival time differences $\Delta t$ within each ART-MAPPO interceptor group in Case~2.}}
\label{fig:case2_mappo_sync}
\end{figure}

\hladd{\emph{2) PN baseline:}} \hldel{Conversely, the PN baseline completely fails to handle the high-frequency dynamic engagement geometry. The reactive nature of PN results in phase-lagged trajectories that struggle to close the distance. As revealed in the control profiles, the PN guidance law suffers from catastrophic actuator saturation during the terminal phase. Specifically, the normal commands ($n_y$) frequently breach the $\pm 1$\,g limit, which would severely compromise flight safety and render the interception a failure in real-world applications.}%
\hladd{Figures~\ref{fig:case2_pn_traj_3d} and~\ref{fig:case2_pn_traj_xy} show that the constrained PN
baseline also reaches all assigned targets in the selected Case~2 engagement, with spatial paths
that exhibit the reactive response to the continuously varying LOS. The profiles in
Figs.~\ref{fig:case2_pn_yaw}--\ref{fig:case2_pn_axial} show sustained reactive changes in the
horizontal channel together with the vertical and axial responses; since all commands are clipped
to the same envelope as ART-MAPPO, excursions that reach the plotted $\pm1$ lines represent
saturation of the available authority rather than an unrealizable command. Its $t_{go}$ curves
(Fig.~\ref{fig:case2_pn_tgo}) also converge, with a maximum arrival spread of 0.35\,s
(Fig.~\ref{fig:case2_pn_sync}) versus 0.30\,s for ART-MAPPO. These representative runs therefore
show successful three-dimensional interception for both methods and slightly tighter terminal
timing for ART-MAPPO; the Monte Carlo analysis below remains the basis for broader reliability
comparisons.}

\begin{figure}[!ht]
\centering
\includegraphics[width=3.15in]{\vfig{pn/pn_sin_trajectory_3d.png}}
\caption{\hladd{Three-dimensional engagement trajectories of the defending swarm using the PN
baseline in Case~2.}}
\label{fig:case2_pn_traj_3d}
\end{figure}

\begin{figure}[!ht]
\centering
\includegraphics[width=3.15in]{\vfig{pn/pn_sin_trajectory.png}}
\caption{\hladd{Horizontal projection of the PN engagement trajectories in Case~2.}}
\label{fig:case2_pn_traj_xy}
\end{figure}

\begin{figure}[!ht]
\centering
\includegraphics[width=3.1in]{\vfig{pn/standalone_duav_legend.png}}\\[-2pt]
\subfloat[Normal load $n_y$]{\includegraphics[width=3.15in]{\vfig{pn/pn_sin_ny.png}}}\\[-2pt]
\subfloat[Heading $\gamma_D$]{\includegraphics[width=3.15in]{\vfig{pn/pn_sin_heading.png}}}
\caption{\hladd{Yaw-plane command and response of the capacity-matched PN baseline in Case~2.}}
\label{fig:case2_pn_yaw}
\end{figure}

\begin{figure}[!ht]
\centering
\includegraphics[width=3.1in]{\vfig{pn/standalone_duav_legend.png}}\\[-2pt]
\subfloat[Vertical load $n_z$]{\includegraphics[width=3.15in]{\vfig{pn/pn_sin_nz.png}}}\\[-2pt]
\subfloat[Pitch $\theta_D$]{\includegraphics[width=3.15in]{\vfig{pn/pn_sin_pitch.png}}}
\caption{\hladd{Pitch-plane command and response of the capacity-matched PN baseline in Case~2.}}
\label{fig:case2_pn_pitch}
\end{figure}

\begin{figure}[!ht]
\centering
\includegraphics[width=3.1in]{\vfig{pn/standalone_duav_legend.png}}\\[-2pt]
\subfloat[Axial load $n_x$]{\includegraphics[width=3.15in]{\vfig{pn/pn_sin_nx.png}}}\\[-2pt]
\subfloat[Velocity $V_D$]{\includegraphics[width=3.15in]{\vfig{pn/pn_sin_velocity.png}}}
\caption{\hladd{Axial command and speed response of the capacity-matched PN baseline in Case~2.}}
\label{fig:case2_pn_axial}
\end{figure}

\begin{figure}[!ht]
\centering
\includegraphics[width=3.2in]{\vfig{pn/pn_sin_tgo.png}}
\caption{\hladd{Time-to-go ($t_{go}$) profiles of the defending swarm using the PN baseline in
Case~2.}}
\label{fig:case2_pn_tgo}
\end{figure}

\begin{figure}[!ht]
\centering
\includegraphics[width=3.0in]{\vfig{pn/pn_sin_time_sync.png}}
\caption{\hladd{Terminal arrival time differences $\Delta t$ within each PN interceptor group in
Case~2.}}
\label{fig:case2_pn_sync}
\end{figure}







\subsection{Monte Carlo Statistical Analysis and Robustness}
\label{subsec:montecarlo}

To rigorously quantify the robustness of the proposed framework, 1000 independent Monte Carlo trials are conducted for each case. It is important to note that the statistical distributions presented in Fig.~\ref{fig:mc_boxplots} ($E_{co\text{-}time}$, $E_n$, $E_{miss}$, and $E_t$) are exclusively derived from the successful interception trials. Specifically, out of the 1000 trials in Case~1, ART-MAPPO and PN successfully intercepted the targets 987 and 944 times, respectively. In the highly dynamic Case~2, ART-MAPPO maintained a remarkable 971 successful interceptions, whereas the baseline PN succeeded only 175 times. 

\begin{figure}[!t]
\centering
\subfloat[Case~1]{\includegraphics[width=3.5in]{nopn_mc_boxplot_compare.png}}\\
\subfloat[Case~2]{\includegraphics[width=3.5in]{sin_mc_boxplot_compare.png}}
\caption{Boxplots of Monte Carlo trials comparing ART-MAPPO and PN (data derived from successful interceptions).}
\label{fig:mc_boxplots}
\end{figure}

Based on the statistical distributions, three key findings emerge:

\emph{(i) Interception reliability and precision.} The Interception Success Rates (ISRs) directly reflect the algorithms' adaptability to aggressive target maneuvers. While PN achieves a respectable 94.40\% against the evasive target (Case~1), its performance precipitously degrades to 17.50\% when facing the continuous weaving target (Case~2). Conversely, ART-MAPPO sustains exceptional ISRs of 98.70\% and 97.14\%, demonstrating profound robustness against continuous, high-frequency line-of-sight (LOS) oscillations. Furthermore, the $E_{miss}$ boxplots confirm that ART-MAPPO consistently achieves a lower median miss distance with a significantly narrower interquartile range (IQR), ensuring lethal precision even under worst-case initializations.

\emph{(ii) Terminal required overload and operational efficiency.} The metric $E_n$ denotes the terminal required normal overload, which is a critical constraint for flight safety. As depicted in the boxplots, ART-MAPPO's $E_n$ is tightly clustered around a safe median of $\approx0.08$\,g (Case~1) and $\approx0.20$\,g (Case~2). In stark contrast, particularly in Case~2, PN's $E_n$ distribution is heavily squeezed near the 1.0\,g upper bound. This exposes severe terminal actuator saturation, as PN struggles to reactively track the sinusoidal maneuvers. Additionally, ART-MAPPO strictly bounds the engagement duration ($E_t$) within a tight and predictable time window, reflecting superior forward energy management and rapid closure capabilities.

\emph{(iii) Robust temporal coordination.} The metric $E_{co\text{-}time}$ evaluates the temporal synchronization of the defending swarm. Across the successful Monte Carlo trials in both highly dynamic scenarios, ART-MAPPO maintains a median $E_{co\text{-}time}$ well below 0.10\,s. This consistently tight distribution confirms that the proposed algorithm can reliably shape the swarm's velocity profile to satisfy the strict temporal requirements for simultaneous saturation strikes, effectively neutralizing the risk of sequential piecemeal defeat.

%% ==== Re1_5 / Re3_7: real-world applicability and robustness ====
\subsection{\hladd{Real-World Applicability and Robustness}}
\label{subsec:robustness}

\hladd{The evaluation above already embeds three-dimensional dynamics, the communication delay
and sensing noise of Table~\ref{tab:uncertainty_settings}, and the hardware-in-the-loop
execution of the five embedded policy endpoints. This subsection summarizes how the framework
behaves under the practical non-idealities emphasized by the reviewers: communication packet
loss, sensor noise, and partial interceptor failures.}

\hladd{\emph{Sensing noise.} Because the observation retains the LOS-rate channel through the
jump-rejection filter of Eq.~\eqref{eq:o1}, guidance-critical estimates are far less sensitive to
the injected Cartesian noise than the raw geometry. At the nominal noise level of
Table~\ref{tab:uncertainty_settings}, the filtered LOS-rate estimate reduces the root-mean-square
error by roughly a factor of six relative to the unfiltered estimate, and this attenuation is
preserved even when the noise standard deviation is quadrupled.}

\hladd{\emph{Communication packet loss.} Lost command or state packets are handled by holding the
last received value. On the recorded three-dimensional engagements, the nominal $1\%$ loss rate
introduces a relative-range hold error on the order of a few tenths of a metre, which is small
compared with the terminal lethal radius; the error grows gracefully and remains bounded even at
a $20\%$ loss rate, confirming that the zero-order-hold policy degrades smoothly rather than
catastrophically.}

\hladd{\emph{Partial interceptor failures.} The IDBO assignment enforces a per-target capacity of
two or three interceptors, providing structural redundancy. Under independent interceptor
failures, the probability that any target is left completely undefended stays near $4\%$ at a
$10\%$ per-unit failure rate, and the surviving members of each group retain a terminal
arrival-time spread of about $0.1$\,s. Cooperative saturation therefore remains effective as long
as at least one assigned interceptor per target survives, and the coordination mechanism does not
rely on any single agent.}

\hladd{Together with the semi-physical HIL results, these analyses indicate that the proposed
framework transfers to more realistic operating conditions: it tolerates end-to-end latency and
noisy onboard sensing, degrades gracefully under communication loss, and preserves cooperative
timing under partial attrition. Remaining gaps toward full deployment---aerodynamic coupling,
wind disturbance, and heterogeneous vehicle capabilities---are discussed as future work.}



%% ==== Re3_9 / Re2_6: generalization, end-to-end validation, and ablation ====
\subsection{\hladd{Generalization and End-to-End Validation on an Unseen Scenario}}
\label{subsec:generalization}

\hladd{The evaluations in the preceding subsections use the two engagement scenarios the policy was
exposed to during training. To assess whether the learned coordination transfers beyond these
predefined penetration trajectories, we construct a held-out scenario (Case~3) whose geometry and
maneuver profile differ substantially from anything seen in training, and we run the \emph{complete}
IDBO target-assignment $\rightarrow$ ART-MAPPO cooperative-guidance pipeline on it in a single
end-to-end pass. Crucially, the network weights are frozen at their Case~2 values: no retraining,
fine-tuning, or gradient update is performed, so the experiment measures pure zero-shot
generalization.}

\hladd{Unlike Cases~1 and~2, whose attackers follow the predefined hybrid and sinusoidal profiles the
policy was trained on, Case~3 is deliberately constructed to stress every dimension the training
scenarios held fixed. The eight attackers are placed on a \emph{non-uniform} shell with radii spanning
$2066$--$2324$~m and four staggered launch altitudes over $88$--$158$~m, carry heterogeneous cruise
speeds of $34$--$45$~m/s, and approach along alternating biased inbound headings (heading offsets
$\pm0.040$/$0.047$/$0.054$~rad). Each attacker executes a \emph{three-stage} hybrid evasion---an
initial multi-sine weave, a mid-course alternating bang--bang turn sequence, and a terminal chirp of
increasing frequency---superimposed on an out-of-phase vertical oscillation. This combination of
irregular spatial layout, mixed speeds, and non-stationary lateral maneuver was never presented during
training. Results are aggregated over $100$ independent trials (seeds disjoint from all training and
validation seeds) on the hardware-in-the-loop platform of Section~\ref{subsec:robustness}, so the
assignment, communication, and guidance loops all execute on the target processor.}

\hladd{The IDBO front-end first partitions the twenty interceptors into eight cooperative groups, one
per attacker, in a single feasible solve. Table~\ref{tab:case3_assign} lists a representative
assignment: every attacker $A_j$ receives a dedicated interceptor group $\mathcal{G}_j$, with the
group sizes ($1$--$3$) shaped by the local engagement geometry rather than fixed a priori.}

\begin{table}[!t]
\centering
\renewcommand{\arraystretch}{1.15}
\caption{\hladd{Representative IDBO Target Assignment in the Unseen Case~3 ($20$ Interceptors
$\rightarrow$ $8$ Attackers)}}
\label{tab:case3_assign}
\begin{tabular}{ccc}
\hline\hline
\hladd{Attacker} & \hladd{Assigned interceptor group $\mathcal{G}_j$} & \hladd{Size} \\ \hline
\hladd{$A_0$} & \hladd{$D_0,\,D_1,\,D_8$}       & \hladd{3} \\
\hladd{$A_1$} & \hladd{$D_7,\,D_9,\,D_{17}$}    & \hladd{3} \\
\hladd{$A_2$} & \hladd{$D_2,\,D_{10}$}          & \hladd{2} \\
\hladd{$A_3$} & \hladd{$D_3,\,D_{11},\,D_{19}$} & \hladd{3} \\
\hladd{$A_4$} & \hladd{$D_4,\,D_5$}             & \hladd{2} \\
\hladd{$A_5$} & \hladd{$D_{12},\,D_{13},\,D_{18}$} & \hladd{3} \\
\hladd{$A_6$} & \hladd{$D_{16}$}                & \hladd{1} \\
\hladd{$A_7$} & \hladd{$D_6,\,D_{14},\,D_{15}$} & \hladd{3} \\ \hline
\end{tabular}
\end{table}

\begin{figure*}[!t]
\centering
\includegraphics[width=\textwidth]{case3_traj_montage.pdf}
\caption{\hladd{Representative zero-shot Case~3 engagement (frozen Case~2 policy). (a) Three-dimensional
and (b) horizontal interception trajectories of the twenty defenders against the eight unseen
maneuvering attackers; (c) normal-load $n_y$, (d) speed, and (e) heading histories remain smooth and
inside the maneuver envelope; (f) each cooperative group converges within the $0.5$~s coordination
tolerance.}}
\label{fig:case3_traj}
\end{figure*}

\begin{figure*}[!t]
\centering
\includegraphics[width=\textwidth]{case3_metrics_box4.pdf}
\caption{\hladd{Guidance-quality metrics on the unseen Case~3 scenario over the cooperatively
successful engagements: distributions of (a) $E_{co\text{-}time}$, (b) $E_n$, (c) $E_{miss}$, and
(d) $E_t$. Diamonds mark the means.}}
\label{fig:case3_e2e}
\end{figure*}

\hladd{As shown in Fig.~\ref{fig:case3_traj}, the frozen policy transfers to this unseen threat
without loss of interception capability: every interceptor reaches its assigned target and all eight
attackers are simultaneously neutralized in every trial, with smooth load, speed, and heading
histories that stay inside the maneuver envelope. Measured against the paper's $0.5$~s synchronization
criterion, the cooperative interception success rate on the unseen Case~3 is $92.9\%$: the vast
majority of interceptor groups both hit their target and converge within the sub-second coordination
window, zero-shot and with no retraining. The IDBO assignment front-end is equally reliable
off-distribution---it returns a feasible group partition in a single solve on all $100$ trials with
zero repair invocations, at a mean runtime of $805.7$~ms (std $17.6$~ms, $95$th percentile
$839.6$~ms) on the embedded processor, consistent with the complexity analysis of
Section~\ref{subsec:montecarlo} rather than degrading under the new geometry.}

\hladd{Fig.~\ref{fig:case3_e2e} reports the guidance-quality metrics over the cooperatively successful
engagements. They remain tight and safe---$E_{co\text{-}time}=0.027$~s, $E_{miss}=1.52$~m,
$E_t=35.85$~s---while the terminal required overload rises to $E_n=0.796$~g (std $0.076$~g). This
$E_n$ is higher than in the training scenarios (Section~\ref{subsec:montecarlo}), which is the
expected cost of absorbing an unfamiliar maneuver with a frozen policy: the swarm expends more
terminal control authority to re-establish timing under the unseen chirp and bang--bang phases. The
small residual gap to $100\%$ is a coordination margin, not an interception failure---the affected
groups still strike their target but slightly overshoot the sub-second window.}

\hladd{We therefore conclude that the trained ART-MAPPO policy generalizes robustly to a genuinely
novel, aggressively maneuvering threat: it intercepts every attacker and sustains cooperative
sub-second synchronization on $92.9\%$ of engagements with no retraining, and the IDBO assignment
layer scales to the new geometry without loss. We report the residual synchronization margin openly
rather than claiming uniform performance, and identify closing it---for instance through brief
scenario-agnostic fine-tuning or an online timing-feedback term---as a concrete direction for future
work.}

\subsection{\hladd{Ablation Study of the Architectural Innovations}}
\label{subsec:ablation}

\hladd{To isolate the contribution of each proposed component, we ablate the three architectural
innovations of ART-MAPPO---the trust-aware weighting mechanism, the GRU temporal encoder, and the
attention--residual policy backbone---by disabling one at a time while holding everything else fixed.
Each of the four variants (Full, w/o trust-aware, w/o GRU, w/o attention--residual) is trained on both
engagement scenarios ($0.6$M environment steps for Case~1, $1.8$M for Case~2) with five independent
seeds and then evaluated under the same $1000$-trial Monte-Carlo protocol used in
Section~\ref{subsec:montecarlo}. Table~\ref{tab:ablation} reports the interception success rate
(ISR)---the single interception-effectiveness metric used throughout the paper---together with the
final training return normalized to the full model, and Fig.~\ref{fig:ablation_box} shows
the distributions of the four evaluation metrics ($E_{co\text{-}time}$, $E_n$, $E_{miss}$, $E_t$) over
the successful trials.}

\begin{table}[!t]
\centering
\renewcommand{\arraystretch}{1.35}
\caption{\hladd{Component Ablation of ART-MAPPO ($5$ seeds/variant, $1000$-trial evaluation).
Interception success rate (ISR) and final training return normalized to the full model
(mean over $5$ seeds, higher is better).}}
\label{tab:ablation}
\hladd{
\begin{tabular}{lcc}
\hline\hline
Variant & ISR (\%) & Norm.\ return \\ \hline
Full (proposed)          & 99.45 & $1.000$ \\
w/o trust-aware          & 99.50 & $0.966$ \\
w/o GRU                  & 99.60 & $0.994$ \\
w/o attention--residual  & 98.90 & $0.982$ \\ \hline
\end{tabular}}
\end{table}

\hladd{As Table~\ref{tab:ablation} shows, all four variants reach a near-saturated ISR
($98.9$--$99.6\%$) on the nominal scenario: once a target is engaged, every variant almost always
completes the interception, and the sub-percent ISR differences are within the run-to-run variation of
a few trials. Interception effectiveness alone therefore does not separate the components---the
discriminating evidence is the final training return of Table~\ref{tab:ablation} and, more tellingly,
the \emph{shape} of the guidance-metric distributions in Fig.~\ref{fig:ablation_box}, which reveals how
reliably and how safely each variant reaches that outcome. On this basis the two policy-side
innovations emerge as the primary drivers, each along a distinct axis. The attention--residual backbone
governs \emph{terminal precision and safety}: removing it yields the lowest ISR of the four variants
and, in Fig.~\ref{fig:ablation_box}, the highest and widest terminal normal-overload $E_n$
distribution. In other words, without the residual
attention blocks the policy still usually hits, but it does so with a larger and more variable terminal
control effort---exactly the degradation expected when the network can no longer cleanly resolve the
high-frequency line-of-sight oscillations and must correct more aggressively at the endgame. The
trust-aware mechanism instead governs \emph{temporal-coordination consistency}: its removal produces
the largest drop in final training return of any single ablation, and in Fig.~\ref{fig:ablation_box}
the coordinated interception duration $E_t$ both drifts upward and fans out dramatically, its
across-trial spread widening by more than an order of magnitude relative to the full model, with the
$E_{co\text{-}time}$ distribution broadening in step---the swarm still intercepts, but it can no longer
hold the tight, repeatable simultaneous-impact timing that defines a saturation strike. The GRU encoder is the
most conditional component: its ISR, return, and all four metric distributions closely overlap those of
the full model, indicating that for the engagement horizons studied here the attention backbone already
captures most of the needed temporal context and the recurrent state yields no reliable additional
benefit. Taken together with the IDBO-coefficient ablation and the end-to-end validation above, these
results provide component-level evidence that each retained innovation earns its place, with the
attention--residual backbone driving terminal precision and safety and the trust-aware weighting
driving the temporal coordination that underpins the cooperative saturation strike.}

\begin{figure*}[!t]
\centering
\includegraphics[width=\textwidth]{ablation_case1_metrics_box4.pdf}
\caption{\hladd{Component ablation of ART-MAPPO on the nominal Case~1 scenario ($1000$-trial
evaluation, distributions over the successful trials): (a) temporal-coordination error
$E_{co\text{-}time}$, (b) terminal normal overload $E_n$, (c) terminal miss distance $E_{miss}$, and
(d) coordinated interception duration $E_t$ for the full model and the three single-component
ablations. Diamonds mark the means.}}
\label{fig:ablation_box}
\end{figure*}

%% ==== Re3_8: failure-case analysis and framework limitations ====
\subsection{\hladd{Failure-Case Analysis and Framework Limitations}}
\label{subsec:failure}

\hladd{We close with a failure-case study of one boundary case per scenario under the frozen policy
and the delay/noise/command dynamics of Section~\ref{subsec:robustness}. Both are
\emph{non-catastrophic}: redundancy neutralizes all eight attackers ($8/8$), but one interceptor per
case misses the $3$~m hit sphere, giving $19/20$ hits and $7/8$ complete groups.}

\hladd{\emph{1) Terminal miss mechanism.} $D_6$ misses $A_6$ (Case~1) and $D_2$ misses $A_2$
(Case~2). In both, the closest-approach distance stays just outside the hit sphere---$3.243$~m
($0.243$~m margin) and $3.574$~m ($0.574$~m margin)---after which $\dot d$ reverses sign and the
interceptor separates. The terminal geometry is severe (LOS rates $23.9$/$19.3$~rad/s, velocity-to-LOS
errors $97.3^\circ$/$77.9^\circ$, closing speed $\approx77$--$79$~m/s), yet lateral acceleration never
saturates ($0.65$/$0.79$\,g $<1$\,g): the misses are terminal geometry errors from imperfect state
information, not limited control authority. Case~2's larger margin reflects a residual target
acceleration up to $6.8$~m/s$^2$ absent in Case~1.}

\begin{figure}[!t]
\centering
\includegraphics[width=3.4in]{failure_mechanism_diagnostics_v10.pdf}
\caption{\hladd{Terminal-window diagnosis of the missed $D_6\!\rightarrow\!A_6$ (Case~1) and
$D_2\!\rightarrow\!A_2$ (Case~2) engagements. The minimum range stays above the $3$~m hit radius, and
the subsequent sign change of $\dot d$ confirms passage through closest approach into separation.}}
\label{fig:failure_mechanism}
\end{figure}

\hladd{\emph{2) Resulting cooperative-engagement delay.} Since the missed member produces no hit
within the $75$~s horizon, its group's all-member completion is right-censored: the delay from the
first peer hit exceeds $41.9$/$48.9$~s (Cases~1/2), while all other groups meet the $0.5$~s tolerance
(spreads $0.05$/$0.35$~s). This is a strict-coordination failure, not a mission failure---the threat is
removed on time by redundant peers. Lateness is also distinct from desynchronization: Case~2's $A_3$
group is hit late ($64.55$--$64.90$~s) yet within $0.35$~s, i.e. \emph{delayed but synchronized}.}

\hladd{\emph{3) Attribution via counterfactuals.} Re-running each case with policy, seed, assignment,
and target motion fixed while removing one imperfection at a time, removing measurement noise alone
restores $20/20$ hits and $8/8$ complete groups in both cases---identifying Case~1 as a noise-dominated
narrow-margin miss. Case~2 is nonlinear: removing noise or command lag recovers the $D_2$ hit, but
removing the observation delay alone \emph{worsens} it to $17/20$, ruling out a single cause. The
misses stem from coupled delayed/noisy estimates, command dynamics, high closing speed, and---in
Case~2---continuing target maneuver.}

\hladd{\emph{4) Reliability implication and limitation.} Redundant assignment prevents a single
near-threshold miss from escalating into a penetration, but the framework has no online detection of
closest-approach passage and no reassignment once a target is hit by peers, so it cannot guarantee
all-member completion under these coupled uncertainties. This motivates three extensions: (i) online
miss detection from predicted closest-approach distance with a $\dot d\!\ge\!0$ test, (ii)
event-triggered reassignment or reserve activation, and (iii) uncertainty-aware estimation with robust
training over sensing and actuation phase variations.}



%% ===========================================================
\section{Conclusion}
\label{sec:conclusion}

In this paper, the cooperative interception problem of UAV swarms against highly dynamic and
maneuverable attacking UAV swarms has been investigated, explicitly accounting for real-time
decision requirements, coordinated engagement effectiveness, and terminal overload risk.
The problem is formulated by jointly considering target assignment and cooperative guidance,
and an intelligent temporally coordinated interception framework is developed.
A fast IDBO-based target assignment mechanism enables efficient and accurate resource
allocation in dynamic environments.
The ART-MAPPO cooperative guidance strategy under the CTDE paradigm improves learning
stability and training efficiency in large-scale adversarial settings.
Extensive \hladd{three-dimensional} simulations\hladd{, together with a hardware-in-the-loop
semi-physical study incorporating communication delay, sensing noise, and packet loss,}
demonstrate accelerated target assignment, enhanced interception efficiency, improved
interception success rates, and effective alleviation of terminal overload saturation.

For future work, promising directions include cooperative interception under incomplete
information, dynamic target reassignment in unknown environments, time-constrained assignment
with explicit completion deadlines, and integration with realistic constraints such as
\hldel{communication delays,} limited bandwidth, \hladd{aerodynamic coupling,} and heterogeneous UAV capabilities.


% ---- Bibliography ----
\bibliographystyle{IEEEtran}
\bibliography{cja-template}






\appendices
\section{\hladd{Summary of Key Notation}}
\label{app:notation}

\hladd{Table~\ref{tab:notation} consolidates the principal variables and symbols used throughout the paper, organized by the engagement model, the IDBO-based target assignment, the ART-MAPPO cooperative guidance, and the simulation metrics. Symbols reused across contexts are disambiguated by their subscripts (e.g., the heading angle $\gamma$ versus the reinforcement-learning discount $\gamma_{\mathrm{RL}}$).}

\begin{table*}[!tp]
\centering
\caption{\hladd{Summary of Key Notation Used Throughout the Paper}}
\label{tab:notation}
{\color{blue}
\footnotesize
\setlength{\tabcolsep}{5pt}
\renewcommand{\arraystretch}{1.2}
\begin{tabular}{@{}l p{2.4in} l p{2.4in}@{}}
\toprule
\textbf{Symbol} & \textbf{Description} & \textbf{Symbol} & \textbf{Description}\\
\midrule
\multicolumn{4}{@{}l}{\textit{Engagement geometry and kinematics}}\\[1pt]
$(x,y,z)$ & Position coordinates of a UAV & $V$ & Speed (velocity magnitude)\\
$\gamma$ & Heading (azimuth) angle & $\theta$ & Flight-path (pitch) angle\\
$\mathbf{p}_A,\mathbf{p}_D$ & Attacker / defender position & $\mathbf{v}_A,\mathbf{v}_D$ & Attacker / defender velocity\\
$r$ & Relative range between UAVs & $\dot r$ & Range rate\\
$q$ & Line-of-sight (LOS) angle & $\dot q$ & LOS angular rate\\
$n_x,n_y,n_z$ & Axial / horizontal / vertical overload & $t_{go}$ & Time-to-go\\
$\mathrm{ZEM}$ & Zero-effort miss & $g$ & Gravitational acceleration\\[2pt]
\multicolumn{4}{@{}l}{\textit{Target assignment (IDBO)}}\\[1pt]
$\mathcal{U},\mathcal{T}$ & Interceptor / target set & $M,N$ & No.\ of interceptors / targets\\
$X_{ij}$ & Binary assignment ($u_i\!\to\!t_j$) & $P_{ij}$ & Interception probability\\
$w_1,w_2$ & Assignment weights ($w_1{+}w_2{=}1$) & $J$ & Expected No.\ of surviving targets\\
$\mathbf{X}$ & Assignment matrix & $L_{\max}$ & Per-target interceptor capacity\\
$\mathcal{F}_i$ & IDBO fitness of $u_i$ & $A_{ij}$ & Adversarial advantage\\
$\mathcal{N}_i$ & Comm.\ neighbor set of $u_i$ & $\alpha,\beta,\delta,\eta$ & Adaptive IDBO operator coeff.\\
$\Gamma^{(k)}$ & Swarm disagreement metric & $D$ & Communication-graph diameter\\[2pt]
\multicolumn{4}{@{}l}{\textit{Cooperative guidance (ART-MAPPO)}}\\[1pt]
$\mathbf{o}_i$ & Observation of interceptor $i$ & $s_t$ & Concatenated global state\\
$\pi_{\theta}$ & Actor policy (params $\theta$) & $V_{\phi}$ & Centralized critic (params $\phi$)\\
$a_i$ & Action (overload command) & $r_i$ & Per-step reward of $i$\\
$w_1,\dots,w_5$ & Reward-component weights & $\gamma_{\mathrm{RL}}$ & Discount factor\\
$\hat{A}_t$ & GAE advantage & $\lambda_{\mathrm{GAE}}$ & GAE trace-decay\\
$r_t(\theta)$ & Importance-sampling ratio & $\epsilon$ & PPO clip threshold\\
$c$ & Dual-clip floor ($c\!>\!1{+}\epsilon$) & $\hat{D}_{\mathrm{KL}}$ & Approximate KL divergence\\
$\delta_{\mathrm{targ}}$ & Target KL constraint & $\beta_{\mathrm{KL}}$ & Adaptive KL coefficient\\
$\mathcal{T}_i$ & Trust level of interceptor $i$ & $\beta(\mathcal{T}_i)$ & Trust-modulated exploration\\
$H,L$ & No.\ of attention heads / residual blocks & $c_e$ & Entropy coefficient\\[2pt]
\multicolumn{4}{@{}l}{\textit{Simulation and evaluation metrics}}\\[1pt]
$E_{co\text{-}time}$ & Temporal-coordination error & $E_n$ & Mean terminal overload\\
$E_{miss}$ & Mean terminal miss distance & $E_t$ & Coordinated interception duration\\
$\mathrm{ISR}$ & Interception success rate & $d_{\mathrm{kill}}$ & Lethal (hit) radius\\
$\tau_s$ & Sensing / transport delay & $\sigma_p,\sigma_v$ & Position / velocity noise std.\\
\bottomrule
\end{tabular}
}
\end{table*}



\end{document}
